% Options for packages loaded elsewhere
\PassOptionsToPackage{unicode}{hyperref}
\PassOptionsToPackage{hyphens}{url}
\documentclass[UTF8,11pt,a4paper]{ctexart}
\usepackage[a4paper,margin=22mm,headheight=15pt]{geometry}
\usepackage{fontspec}
\usepackage{xeCJK}
\setmainfont{Segoe UI}
\setsansfont{Segoe UI}
\setmonofont{Cascadia Mono}
% Noto Sans SC is distributed as a regular face on this Windows image. Fake
% the missing emphasis faces so headings and inline code keep one CJK family.
\setCJKmainfont[AutoFakeBold=2,AutoFakeSlant=.2]{Noto Sans SC}
\setCJKsansfont[AutoFakeBold=2,AutoFakeSlant=.2]{Noto Sans SC}
\setCJKmonofont[AutoFakeBold=2,AutoFakeSlant=.2]{Noto Sans SC}
\usepackage{xcolor}
\usepackage{amsmath,amssymb}
% The Markdown headings already carry the guide's stable section numbers.
% Disable LaTeX's second numbering layer so the TOC stays readable.
\setcounter{secnumdepth}{0}
\usepackage{iftex}
\ifPDFTeX
  \usepackage[T1]{fontenc}
  \usepackage[utf8]{inputenc}
  \usepackage{textcomp} % provide euro and other symbols
\else % if luatex or xetex
  \usepackage{unicode-math} % this also loads fontspec
  \defaultfontfeatures{Scale=MatchLowercase}
  \defaultfontfeatures[\rmfamily]{Ligatures=TeX,Scale=1}
\fi
\usepackage{lmodern}
\ifPDFTeX\else
  % xetex/luatex font selection
\fi
% Use upquote if available, for straight quotes in verbatim environments
\IfFileExists{upquote.sty}{\usepackage{upquote}}{}
\IfFileExists{fvextra.sty}{\usepackage{fvextra}}{}
\IfFileExists{microtype.sty}{% use microtype if available
  \usepackage[]{microtype}
  \UseMicrotypeSet[protrusion]{basicmath} % disable protrusion for tt fonts
}{}
\makeatletter
\@ifundefined{KOMAClassName}{% if non-KOMA class
  \IfFileExists{parskip.sty}{%
    \usepackage{parskip}
  }{% else
    \setlength{\parindent}{0pt}
    \setlength{\parskip}{6pt plus 2pt minus 1pt}}
}{% if KOMA class
  \KOMAoptions{parskip=half}}
\makeatother
\usepackage{color}
\usepackage{fancyvrb}
\newcommand{\VerbBar}{|}
\newcommand{\VERB}{\Verb[commandchars=\\\{\}]}
\newenvironment{Highlighting}{%
  \VerbatimEnvironment%
  \begin{Verbatim}[commandchars=\\\{\},breaklines=true,breakanywhere=true,fontsize=\small,baselinestretch=1.15]}{\end{Verbatim}}
% Add ',fontsize=\small' for more characters per line
\usepackage{framed}
\definecolor{shadecolor}{RGB}{248,248,248}
\newenvironment{Shaded}{\begin{snugshade}}{\end{snugshade}}
\newcommand{\AlertTok}[1]{\textcolor[rgb]{0.94,0.16,0.16}{#1}}
\newcommand{\AnnotationTok}[1]{\textcolor[rgb]{0.56,0.35,0.01}{\textbf{\textit{#1}}}}
\newcommand{\AttributeTok}[1]{\textcolor[rgb]{0.13,0.29,0.53}{#1}}
\newcommand{\BaseNTok}[1]{\textcolor[rgb]{0.00,0.00,0.81}{#1}}
\newcommand{\BuiltInTok}[1]{#1}
\newcommand{\CharTok}[1]{\textcolor[rgb]{0.31,0.60,0.02}{#1}}
\newcommand{\CommentTok}[1]{\textcolor[rgb]{0.56,0.35,0.01}{\textit{#1}}}
\newcommand{\CommentVarTok}[1]{\textcolor[rgb]{0.56,0.35,0.01}{\textbf{\textit{#1}}}}
\newcommand{\ConstantTok}[1]{\textcolor[rgb]{0.56,0.35,0.01}{#1}}
\newcommand{\ControlFlowTok}[1]{\textcolor[rgb]{0.13,0.29,0.53}{\textbf{#1}}}
\newcommand{\DataTypeTok}[1]{\textcolor[rgb]{0.13,0.29,0.53}{#1}}
\newcommand{\DecValTok}[1]{\textcolor[rgb]{0.00,0.00,0.81}{#1}}
\newcommand{\DocumentationTok}[1]{\textcolor[rgb]{0.56,0.35,0.01}{\textbf{\textit{#1}}}}
\newcommand{\ErrorTok}[1]{\textcolor[rgb]{0.64,0.00,0.00}{\textbf{#1}}}
\newcommand{\ExtensionTok}[1]{#1}
\newcommand{\FloatTok}[1]{\textcolor[rgb]{0.00,0.00,0.81}{#1}}
\newcommand{\FunctionTok}[1]{\textcolor[rgb]{0.13,0.29,0.53}{\textbf{#1}}}
\newcommand{\ImportTok}[1]{#1}
\newcommand{\InformationTok}[1]{\textcolor[rgb]{0.56,0.35,0.01}{\textbf{\textit{#1}}}}
\newcommand{\KeywordTok}[1]{\textcolor[rgb]{0.13,0.29,0.53}{\textbf{#1}}}
\newcommand{\NormalTok}[1]{#1}
\newcommand{\OperatorTok}[1]{\textcolor[rgb]{0.81,0.36,0.00}{\textbf{#1}}}
\newcommand{\OtherTok}[1]{\textcolor[rgb]{0.56,0.35,0.01}{#1}}
\newcommand{\PreprocessorTok}[1]{\textcolor[rgb]{0.56,0.35,0.01}{\textit{#1}}}
\newcommand{\RegionMarkerTok}[1]{#1}
\newcommand{\SpecialCharTok}[1]{\textcolor[rgb]{0.81,0.36,0.00}{\textbf{#1}}}
\newcommand{\SpecialStringTok}[1]{\textcolor[rgb]{0.31,0.60,0.02}{#1}}
\newcommand{\StringTok}[1]{\textcolor[rgb]{0.31,0.60,0.02}{#1}}
\newcommand{\VariableTok}[1]{\textcolor[rgb]{0.00,0.00,0.00}{#1}}
\newcommand{\VerbatimStringTok}[1]{\textcolor[rgb]{0.31,0.60,0.02}{#1}}
\newcommand{\WarningTok}[1]{\textcolor[rgb]{0.56,0.35,0.01}{\textbf{\textit{#1}}}}
\usepackage{longtable,booktabs,array}
\usepackage{seqsplit}
\usepackage{colortbl}
\setlength{\tabcolsep}{3pt}
\renewcommand{\arraystretch}{1.24}
\definecolor{tablehead}{RGB}{232,240,252}
\definecolor{tablealt}{RGB}{248,250,253}
\newcolumntype{L}[1]{>{\raggedright\arraybackslash}p{#1}}
\newcommand{\themeclass}[1]{\begingroup\footnotesize\texttt{\seqsplit{#1}}\endgroup}
\usepackage{calc} % for calculating minipage widths
% Correct order of tables after \paragraph or \subparagraph
\usepackage{etoolbox}
\makeatletter
\patchcmd\longtable{\par}{\if@noskipsec\mbox{}\fi\par}{}{}
\makeatother
% Allow footnotes in longtable head/foot
\IfFileExists{footnotehyper.sty}{\usepackage{footnotehyper}}{\usepackage{footnote}}
\makesavenoteenv{longtable}
\setlength{\LTpre}{8pt plus 2pt minus 2pt}
\setlength{\LTpost}{10pt plus 2pt minus 2pt}
\AtBeginEnvironment{longtable}{\rowcolors{2}{tablealt}{white}}
\AtEndEnvironment{longtable}{\hiderowcolors}
\setlength{\emergencystretch}{3em} % prevent overfull lines
\providecommand{\tightlist}{%
  \setlength{\itemsep}{0pt}\setlength{\parskip}{0pt}}
\usepackage{bookmark}
\IfFileExists{xurl.sty}{\usepackage{xurl}}{} % add URL line breaks if available
\urlstyle{same}
\usepackage[most]{tcolorbox}
\usepackage{enumitem}
\usepackage{fancyhdr}
\usepackage{titlesec}
\setlist[itemize]{nosep,leftmargin=*,topsep=2pt,partopsep=0pt}
\setlist[enumerate]{nosep,leftmargin=*,topsep=2pt,partopsep=0pt}
\setlength{\parindent}{0pt}
\setlength{\parskip}{5pt plus 1pt minus 1pt}
\widowpenalty=10000
\clubpenalty=10000
\brokenpenalty=10000
\emergencystretch=4em
\titleformat{\section}{\Large\bfseries\color{blue!55!black}}{\thesection}{0.8em}{}
\titleformat{\subsection}{\large\bfseries\color{blue!45!black}}{\thesubsection}{0.7em}{}
\titleformat{\subsubsection}{\normalsize\bfseries\color{black!75}}{\thesubsubsection}{0.6em}{}
\pagestyle{fancy}
\fancyhf{}
\lhead{FanControl 主题 DIY 完整指南}
\rhead{设计规范与代码说明}
\cfoot{\thepage}
\hypersetup{
  colorlinks=true,
  linkcolor=blue!55!black,
  urlcolor=blue!65!black,
  citecolor=blue!55!black,
  pdftitle={FanControl 主题 DIY 完整指南},
  pdfauthor={FanControl Project}}
\hypersetup{
  pdftitle={FanControl 主题 DIY 完整指南},
  pdfauthor={FanControl Project},
  pdfcreator={LaTeX via pandoc}}

\title{FanControl 主题 DIY 完整指南}
\author{FanControl Project}
\date{2026-07-17}

\begin{document}
\maketitle

{
\setcounter{tocdepth}{3}
\tableofcontents
}
\section{FanControl 主题 DIY
完整指南}\label{fancontrol-ux4e3bux9898-diy-ux5b8cux6574ux6307ux5357}

\begin{quote}
适用于 FanControl 2.5.x
及之后版本的基础主题和高级主题。本文既是设计说明，也是可以交给作者、设计师或
AI 使用的工程规范。
\end{quote}

本文解释主题包的每一个文件、字段、CSS
作用域、页面钩子、组件钩子、资源路径、字体子集、调试方法和发布验收规则。主题只修改
\texttt{theme.json}、\texttt{theme.css} 和主题资源，不修改 FanControl
源码。

\begin{center}\rule{0.5\linewidth}{0.5pt}\end{center}

\subsection{0.
先理解主题系统}\label{0-ux5148ux7406ux89e3ux4e3bux9898ux7cfbux7edf}

FanControl 的主题不是一个独立插件进程，也不是一个网页。它是一组被 GUI
读取的静态资源：

\begin{enumerate}
\def\labelenumi{\arabic{enumi}.}
\tightlist
\item
  GUI 启动时，后端扫描可执行文件旁边的 \texttt{themes/}
  目录，并读取每个主题的 \texttt{theme.json}。
\item
  前端主题同步组件根据配置中的 \texttt{themeMode}
  决定使用系统、浅色、深色或自定义主题。
\item
  选择自定义主题时，GUI 后端读取该主题的 \texttt{theme.css}。
\item
  CSS 被放入页面的
  \texttt{\textless{}style\ id="thrm-custom-theme-style"\textgreater{}}，并给
  \texttt{\textless{}html\textgreater{}} 加上 \texttt{data-theme} 属性。
\item
  CSS 中的 \texttt{html{[}data-theme="主题ID"{]}}
  选择器因此只影响当前主题。
\item
  CSS 引用的本地图片、SVG 和字体通过
  \texttt{/theme-assets/\textless{}主题ID\textgreater{}/\textless{}相对路径\textgreater{}}
  提供。
\end{enumerate}

主题只运行在 GUI 层，不会进入 FanControl
Core，也不会直接触碰设备、温度采集或风扇控制逻辑。联网下载主题时，网络任务同样应该属于
GUI，不应阻塞 Core 的监控循环。

\subsubsection{0.1
主题加载关系}\label{01-ux4e3bux9898ux52a0ux8f7dux5173ux7cfb}

\begin{Shaded}
\begin{Highlighting}
\NormalTok{theme.json}
\NormalTok{    |}
\NormalTok{    v}
\NormalTok{主题列表 {-}{-}{-}{-}{-}\textgreater{} 设置页主题选择器}
\NormalTok{    |}
\NormalTok{    v}
\NormalTok{theme.css {-}{-}{-}{-}{-}\textgreater{} \textless{}style\textgreater{} + html[data{-}theme="id"]}
\NormalTok{    |}
\NormalTok{    v}
\NormalTok{相对资源 {-}{-}{-}{-}{-}\textgreater{} /theme{-}assets/id/path}
\NormalTok{    |}
\NormalTok{    v}
\NormalTok{页面组件上的 glacier{-}* / data{-}theme{-}* 钩子}
\end{Highlighting}
\end{Shaded}

上面的每一条箭头都对应一个约束：清单负责``我是谁''，CSS
负责``怎么画''，资源路由负责``从哪里取文件''，组件钩子负责``影响哪一个界面部件''。

\subsubsection{0.2
当前实现中可以复用的代码}\label{02-ux5f53ux524dux5b9eux73b0ux4e2dux53efux4ee5ux590dux7528ux7684ux4ee3ux7801}

主题作者通常不需要阅读 Go 代码，但理解下面几个入口有助于排查问题：

\begin{longtable}[]{@{}L{0.36\textwidth}L{0.58\textwidth}@{}}
\toprule\noalign{}\rowcolor{tablehead}
文件 & 作用 \\
\midrule\noalign{}
\endhead
\bottomrule\noalign{}
\endlastfoot
\texttt{internal/theme/theme.go} & 扫描清单、迁移旧主题、读取
CSS、读取资源、校验主题 ID 和资源路径 \\
\texttt{internal/guiapp/theme\_api.go} & 向前端提供
\texttt{ListThemes}、\texttt{GetThemeCSS}、\texttt{OpenThemesFolder} \\
\texttt{main.go} & 注册 \texttt{/theme-assets/} 资源处理器 \\
\shortstack[l]{\texttt{frontend/src/app/}\\
  \texttt{components/}\\
  \texttt{SystemThemeSync.tsx}} &
读取配置并应用系统或自定义主题 \\
\shortstack[l]{\texttt{frontend/src/app/}\\
  \texttt{components/settings/}\\
  \texttt{SystemSettingsSection.tsx}}
& 在设置页加载主题列表并保存选择 \\
\texttt{frontend/src/app/globals.css} & 定义 FanControl
本身的基础外壳和高级钩子 \\
\texttt{docs/theme-template/basic/} & 兼容性优先的基础模板 \\
\texttt{docs/theme-template/advanced/} &
可使用页面、卡片、字体、贴纸和背景资源的高级模板 \\
\end{longtable}

\begin{center}\rule{0.5\linewidth}{0.5pt}\end{center}

\subsection{1. 选择模板：Basic 还是
Advanced}\label{1-ux9009ux62e9ux6a21ux677fbasic-ux8fd8ux662f-advanced}

\subsubsection{1.1 Basic
基础主题}\label{11-basic-ux57faux7840ux4e3bux9898}

Basic 适合颜色、背景、卡片、图表和左侧徽标调整。它只依赖稳定的 CSS
变量、\texttt{.glacier-*} 外壳类和原生 HTML 属性，适合希望同时兼容 THRM
或旧版本的作者。

Basic 的设计目标是``少改结构、保持可读''。它不应该依赖 FanControl 新增的
\texttt{data-theme-card} 和 \texttt{data-theme-ui}
细粒度钩子。这样即使旧版本忽略某些高级选择器，基础颜色和布局仍然可用。

\subsubsection{1.2 Advanced
高级主题}\label{12-advanced-ux9ad8ux7ea7ux4e3bux9898}

Advanced
适合完整视觉设计，包括首页插画、背景幕布、贴纸、字体子集、页面级背景、卡片透明度、hover
动效和高级组件钩子。它主要面向 FanControl 2.3.0 及之后版本。

Advanced
仍然应该保留基础层的颜色变量和可读性兜底。高级选择器被旧软件忽略时，页面应该退化为清晰的基础主题，而不是变成无法操作的空白界面。

\subsubsection{1.3 选择决策}\label{13-ux9009ux62e9ux51b3ux7b56}

\begin{longtable}[]{@{}L{0.36\textwidth}L{0.58\textwidth}@{}}
\toprule\noalign{}\rowcolor{tablehead}
需求 & 推荐 \\
\midrule\noalign{}
\endhead
\bottomrule\noalign{}
\endlastfoot
只改颜色、圆角、卡片和图表 & Basic \\
要兼容 THRM 或旧版本 & Basic \\
要使用首页插画、幕布和贴纸 & Advanced \\
要绑定 FanControl 的首页、曲线页、设置页 & Advanced \\
要使用主题专用字体 & Advanced，并使用字体子集 \\
不确定自己能否维护复杂 CSS & Basic \\
\end{longtable}

\begin{center}\rule{0.5\linewidth}{0.5pt}\end{center}

\subsection{2.
主题目录结构}\label{2-ux4e3bux9898ux76eeux5f55ux7ed3ux6784}

\subsubsection{2.1
最小可用结构}\label{21-ux6700ux5c0fux53efux7528ux7ed3ux6784}

\begin{Shaded}
\begin{Highlighting}
\NormalTok{my{-}theme/}
\NormalTok{  theme.json       \# 主题身份和元数据}
\NormalTok{  theme.css        \# 主题样式}
\NormalTok{  README.md        \# 可选：作者说明、许可证、预览说明}
\end{Highlighting}
\end{Shaded}

\subsubsection{2.2
高级主题结构}\label{22-ux9ad8ux7ea7ux4e3bux9898ux7ed3ux6784}

\begin{Shaded}
\begin{Highlighting}
\NormalTok{my{-}theme/}
\NormalTok{  theme.json}
\NormalTok{  theme.css}
\NormalTok{  README.md}
\NormalTok{  hero.webp                    \# 首页插画或主视觉}
\NormalTok{  logo.svg                     \# 可选：侧栏徽标}
\NormalTok{  assets/}
\NormalTok{    fonts/}
\NormalTok{      my{-}theme{-}ui.woff2       \# 可选：UI 字体子集}
\NormalTok{      LICENSE{-}font.txt        \# 字体许可证}
\NormalTok{    stickers/}
\NormalTok{      curtain.webp             \# 可选：背景幕布}
\NormalTok{  decorations/}
\NormalTok{    cloud.svg}
\NormalTok{    star.svg}
\end{Highlighting}
\end{Shaded}

\subsubsection{2.3 分享和上传时的 ZIP
结构}\label{23-ux5206ux4eabux548cux4e0aux4f20ux65f6ux7684-zip-ux7ed3ux6784}

ZIP 内建议保留一层主题目录：

\begin{Shaded}
\begin{Highlighting}
\NormalTok{my{-}theme{-}1.0.0.zip}
\NormalTok{  my{-}theme/}
\NormalTok{    theme.json}
\NormalTok{    theme.css}
\NormalTok{    README.md}
\NormalTok{    hero.webp}
\NormalTok{    assets/}
\NormalTok{    decorations/}
\end{Highlighting}
\end{Shaded}

这样用户解压后可以直接把 \texttt{my-theme/} 放进
\texttt{themes/}，也便于安装流程进行目录校验。

不要放入主题包：

\begin{longtable}[]{@{}L{0.36\textwidth}L{0.58\textwidth}@{}}
\toprule\noalign{}\rowcolor{tablehead}
不应包含 & 原因 \\
\midrule\noalign{}
\endhead
\bottomrule\noalign{}
\endlastfoot
\texttt{Cache/}、\texttt{node\_modules/}、\texttt{.git/} &
开发过程文件，会增大包体积 \\
PSD、AI、工程源文件 & 设计源文件不属于运行时资源 \\
截图原图和中间渲染图 & 只保留最终预览和运行时资源 \\
\texttt{C:\textbackslash{}Users\textbackslash{}...}、\texttt{D:\textbackslash{}...}
路径 & 其他电脑无法读取 \\
\texttt{.exe}、\texttt{.dll}、\texttt{.js}、\texttt{.html} & 主题只需要
CSS 和静态资源 \\
未授权的完整商业字体 & 可能违反字体许可证 \\
\end{longtable}

\begin{center}\rule{0.5\linewidth}{0.5pt}\end{center}

\subsection{3.
第一次创建主题}\label{3-ux7b2cux4e00ux6b21ux521bux5efaux4e3bux9898}

\subsubsection{3.1 复制模板}\label{31-ux590dux5236ux6a21ux677f}

基础主题从这里复制：

\begin{Shaded}
\begin{Highlighting}
\NormalTok{docs/theme{-}template/basic/}
\end{Highlighting}
\end{Shaded}

高级主题从这里复制：

\begin{Shaded}
\begin{Highlighting}
\NormalTok{docs/theme{-}template/advanced/}
\end{Highlighting}
\end{Shaded}

不要直接在内置主题目录上改动。复制模板的目的，是保留主题所需的变量兜底、页面结构兼容和注释。

\subsubsection{3.2 修改主题 ID}\label{32-ux4feeux6539ux4e3bux9898-id}

假设新主题叫 \texttt{cute-blue}，需要同步修改三个位置：

\begin{enumerate}
\def\labelenumi{\arabic{enumi}.}
\tightlist
\item
  文件夹名改为 \texttt{cute-blue}。
\item
  \texttt{theme.json} 中的 \texttt{id} 改为 \texttt{cute-blue}。
\item
  \texttt{theme.css} 中所有 \texttt{html{[}data-theme="advanced"{]}} 或
  \texttt{html{[}data-theme="basic"{]}} 改为
  \texttt{html{[}data-theme="cute-blue"{]}}。
\end{enumerate}

如果只改了文件夹或只改了 JSON，主题可能出现在列表里，但 CSS 不会生效。ID
同步是最常见、也最容易漏掉的步骤。

\subsubsection{3.3 ID 命名规则}\label{33-id-ux547dux540dux89c4ux5219}

主题 ID 必须满足：

\begin{Shaded}
\begin{Highlighting}
\NormalTok{\^{}[a{-}z0{-}9\_{-}]\{1,64\}$}
\end{Highlighting}
\end{Shaded}

含义如下：

\begin{longtable}[]{@{}L{0.36\textwidth}L{0.58\textwidth}@{}}
\toprule\noalign{}\rowcolor{tablehead}
片段 & 意思 \\
\midrule\noalign{}
\endhead
\bottomrule\noalign{}
\endlastfoot
\texttt{a-z} & 只能使用小写英文字母 \\
\texttt{0-9} & 可以使用数字 \\
\texttt{-} & 短横线可用于分词 \\
\texttt{\_} & 下划线可用于兼容旧命名 \\
\texttt{\{1,64\}} & 长度至少 1，最多 64 \\
\end{longtable}

显示名称 \texttt{name} 不受这个限制，可以写中文、空格和品牌名称。

\begin{center}\rule{0.5\linewidth}{0.5pt}\end{center}

\subsection{\texorpdfstring{4. \texttt{theme.json}
逐字段说明}{4. theme.json 逐字段说明}}\label{4-themejson-ux9010ux5b57ux6bb5ux8bf4ux660e}

JSON 不支持 \texttt{//} 注释。模板使用 \texttt{\$comment} 和
\texttt{\$help}，因为它们是普通 JSON 字段，解析器可以安全忽略。

\subsubsection{4.1 完整示例}\label{41-ux5b8cux6574ux793aux4f8b}

\begin{Shaded}
\begin{Highlighting}
\FunctionTok{\{}
  \DataTypeTok{"$comment"}\FunctionTok{:} \StringTok{"这里写给主题作者看的说明，运行时会忽略。"}\FunctionTok{,}
  \DataTypeTok{"$help"}\FunctionTok{:} \FunctionTok{\{}
    \DataTypeTok{"id"}\FunctionTok{:} \StringTok{"唯一 ID，必须和文件夹名一致。"}\FunctionTok{,}
    \DataTypeTok{"name"}\FunctionTok{:} \StringTok{"显示在设置页里的名称。"}\FunctionTok{,}
    \DataTypeTok{"base"}\FunctionTok{:} \StringTok{"light 或 dark。"}\FunctionTok{,}
    \DataTypeTok{"layer"}\FunctionTok{:} \StringTok{"basic 或 advanced。"}\FunctionTok{,}
    \DataTypeTok{"author"}\FunctionTok{:} \StringTok{"作者名。"}\FunctionTok{,}
    \DataTypeTok{"version"}\FunctionTok{:} \StringTok{"主题版本。"}\FunctionTok{,}
    \DataTypeTok{"description"}\FunctionTok{:} \StringTok{"主题简介。"}
  \FunctionTok{\},}
  \DataTypeTok{"id"}\FunctionTok{:} \StringTok{"cute{-}blue"}\FunctionTok{,}
  \DataTypeTok{"name"}\FunctionTok{:} \StringTok{"Cute Blue"}\FunctionTok{,}
  \DataTypeTok{"base"}\FunctionTok{:} \StringTok{"light"}\FunctionTok{,}
  \DataTypeTok{"layer"}\FunctionTok{:} \StringTok{"advanced"}\FunctionTok{,}
  \DataTypeTok{"author"}\FunctionTok{:} \StringTok{"FanControl User"}\FunctionTok{,}
  \DataTypeTok{"version"}\FunctionTok{:} \StringTok{"1.0.0"}\FunctionTok{,}
  \DataTypeTok{"description"}\FunctionTok{:} \StringTok{"蓝白配色的轻量高级主题。"}
\FunctionTok{\}}
\end{Highlighting}
\end{Shaded}

\subsubsection{4.2 字段解释}\label{42-ux5b57ux6bb5ux89e3ux91ca}

\begin{longtable}[]{@{}L{0.14\textwidth}L{0.10\textwidth}L{0.25\textwidth}L{0.41\textwidth}@{}}
\toprule\noalign{}\rowcolor{tablehead}
字段 & 必填 & 允许值 & 运行时作用 \\
\midrule\noalign{}
\endhead
\bottomrule\noalign{}
\endlastfoot
\texttt{\$comment} & 否 & 任意字符串 & 给作者阅读，不参与主题身份 \\
\texttt{\$help} & 否 & 对象 & 给复制模板的人阅读，不参与主题身份 \\
\texttt{id} & 推荐 & 小写字母、数字、\texttt{-}、\texttt{\_} &
目录名、\texttt{data-theme} 值和资源路由中的主题名 \\
\texttt{name} & 否 & 任意字符串 & 设置页显示名称；空时退回 ID \\
\texttt{base} & 否 & \texttt{light} 或 \texttt{dark} & 应用主题 CSS
前使用的明暗基础 \\
\texttt{layer} & 否 & \texttt{basic} 或 \texttt{advanced} &
告诉前端该主题是否使用高级钩子 \\
\texttt{interface} & 兼容 & \texttt{basic} 或 \texttt{advanced} &
老模板字段，会被当作 \texttt{layer} 使用 \\
\texttt{author} & 否 & 任意字符串 & 设置页和主题说明中显示作者 \\
\texttt{version} & 否 & 建议使用 \texttt{主.次.修} &
主题升级和迁移时比较版本 \\
\texttt{description} & 否 & 任意字符串 & 主题列表和作者说明 \\
\end{longtable}

\subsubsection{\texorpdfstring{4.3 \texttt{base}
的含义}{4.3 base 的含义}}\label{43-base-ux7684ux542bux4e49}

\texttt{base}
不是主题的完整样式，它只是告诉应用先使用浅色还是深色基础变量：

\begin{Shaded}
\begin{Highlighting}
\NormalTok{base = light  {-}\textgreater{} html 不加 dark 类}
\NormalTok{base = dark   {-}\textgreater{} html 加 dark 类}
\end{Highlighting}
\end{Shaded}

自定义主题不是 \texttt{system}。如果用户选择自定义主题，主题自己的
\texttt{base} 决定明暗；系统深浅变化不会自动把自定义主题改成另一套颜色。

\subsubsection{\texorpdfstring{4.4 \texttt{layer}
的含义}{4.4 layer 的含义}}\label{44-layer-ux7684ux542bux4e49}

\texttt{basic} 代表只依赖通用结构；\texttt{advanced} 代表可以使用
FanControl 高级钩子。该字段不会自动授予主题权限，也不会替作者生成
CSS，只用于前端选择正确的兼容层和主题状态。

\subsubsection{4.5 版本号规则}\label{45-ux7248ux672cux53f7ux89c4ux5219}

建议使用语义化版本：

\begin{Shaded}
\begin{Highlighting}
\NormalTok{1.0.0       首个公开版本}
\NormalTok{1.1.0       新增视觉功能或资源}
\NormalTok{1.1.1       修复颜色、路径或可读性问题}
\NormalTok{2.0.0       大幅改变布局或资源结构}
\end{Highlighting}
\end{Shaded}

不要用时间戳代替版本号。安装器会按数字片段比较版本；\texttt{1.10.0}
应该高于 \texttt{1.9.0}。

\begin{center}\rule{0.5\linewidth}{0.5pt}\end{center}

\subsection{5. CSS
的核心规则：作用域、顺序和变量}\label{5-css-ux7684ux6838ux5fc3ux89c4ux5219ux4f5cux7528ux57dfux987aux5e8fux548cux53d8ux91cf}

\subsubsection{5.1
所有选择器都要进入主题作用域}\label{51-ux6240ux6709ux9009ux62e9ux5668ux90fdux8981ux8fdbux5165ux4e3bux9898ux4f5cux7528ux57df}

正确写法：

\begin{Shaded}
\begin{Highlighting}
\CommentTok{/* 主题入口：下面所有规则只对 cute{-}blue 生效。 */}
\NormalTok{html}\ExtensionTok{[}\SpecialStringTok{data{-}theme}\OperatorTok{=}\StringTok{"cute{-}blue"}\ExtensionTok{]}\NormalTok{ \{}
  \VariableTok{{-}{-}primary}\CharTok{:} \ConstantTok{\#2f6df6}\OperatorTok{;} \CommentTok{/* 主色，按钮、选中态和强调文字会复用。 */}
\NormalTok{\}}

\CommentTok{/* 只修改主题外壳，不影响其他主题。 */}
\NormalTok{html}\ExtensionTok{[}\SpecialStringTok{data{-}theme}\OperatorTok{=}\StringTok{"cute{-}blue"}\ExtensionTok{]} \FunctionTok{.glacier{-}shell}\NormalTok{ \{}
  \KeywordTok{background}\CharTok{:} \FunctionTok{var(}\VariableTok{{-}{-}background}\FunctionTok{)}\OperatorTok{;} \CommentTok{/* 使用应用已经计算好的背景变量。 */}
\NormalTok{\}}
\end{Highlighting}
\end{Shaded}

错误写法：

\begin{Shaded}
\begin{Highlighting}
\CommentTok{/* 错误：会影响系统主题和其他自定义主题。 */}
\NormalTok{body \{}
  \KeywordTok{background}\CharTok{:} \ConstantTok{red}\OperatorTok{;}
\NormalTok{\}}

\CommentTok{/* 错误：没有 data{-}theme 限制作用域。 */}
\FunctionTok{.glacier{-}shell}\NormalTok{ \{}
  \KeywordTok{background}\CharTok{:} \ConstantTok{red}\OperatorTok{;}
\NormalTok{\}}
\end{Highlighting}
\end{Shaded}

\subsubsection{\texorpdfstring{5.2 为什么不应该直接覆盖
\texttt{body}}{5.2 为什么不应该直接覆盖 body}}\label{52-ux4e3aux4ec0ux4e48ux4e0dux5e94ux8be5ux76f4ux63a5ux8986ux76d6-body}

FanControl 的窗口背景、原生材质、系统明暗和自定义主题是多层组合。全局
\texttt{body}
规则容易覆盖基础主题的透明处理，也可能让切换回系统主题后留下自定义颜色。主题应当从
\texttt{html{[}data-theme="id"{]}} 开始，并尽量改变量而不是重写布局。

\subsubsection{5.3 CSS
变量覆盖优先于重复布局}\label{53-css-ux53d8ux91cfux8986ux76d6ux4f18ux5148ux4e8eux91cdux590dux5e03ux5c40}

推荐顺序：

\begin{enumerate}
\def\labelenumi{\arabic{enumi}.}
\tightlist
\item
  先修改现有 CSS 变量。
\item
  再修改 \texttt{.glacier-*} 的颜色、边框、阴影和背景。
\item
  只有确实需要时，才使用 \texttt{data-theme-card} 和
  \texttt{data-theme-ui}。
\item
  不要通过负边距、绝对定位和大范围 \texttt{!important} 重建整个页面。
\end{enumerate}

\subsubsection{5.4 CSS
文件组织顺序}\label{54-css-ux6587ux4ef6ux7ec4ux7ec7ux987aux5e8f}

建议把 \texttt{theme.css} 分为以下区块，并用注释标明：

\begin{Shaded}
\begin{Highlighting}
\CommentTok{/* 1. Theme identity: ID prefix and main variables. */}
\CommentTok{/* 2. Global shell: window, title bar, sidebar, content. */}
\CommentTok{/* 3. Cards and controls: hero, metrics, charts, inputs. */}
\CommentTok{/* 4. Page hooks: status, curve, control, about, devices. */}
\CommentTok{/* 5. Assets and fonts: @font{-}face and local URLs. */}
\CommentTok{/* 6. Responsive rules: narrow windows and reduced motion. */}
\CommentTok{/* 7. Final guards: clipping, readability, and known component fixes. */}
\end{Highlighting}
\end{Shaded}

文件末尾的``Final
guards''不是鼓励堆叠覆盖，而是给已知组件留下一个清晰的最终校正区。每条覆盖都应说明它修复了哪个组件和为什么放在末尾。

\subsubsection{5.5 代码块的详细规范}

代码块不是把几行 CSS 粘贴进文档就结束了。每个示例都应该让读者能回答三个问题：它作用于哪个节点、它解决什么视觉问题、删掉它会发生什么。推荐在代码块前写一句目的说明，在复杂规则旁保留一条“原因 + 约束”注释。

\begin{longtable}[]{@{}L{0.20\textwidth}L{0.34\textwidth}L{0.40\textwidth}@{}}
\toprule\noalign{}\rowcolor{tablehead}
代码位置 & 必须说明 & 示例问题 \\
\midrule\noalign{}
\endhead
\bottomrule\noalign{}
\endlastfoot
选择器第一行 & 页面、卡片或控件范围 & 这是首页主卡片，还是所有页面的按钮？ \\
CSS 变量区 & 变量的语义和回退 & \texttt{--cute-blue-border} 是边框色，不是文字色 \\
视觉属性 & 修改后影响的层级 & \texttt{box-shadow} 只表达层级，不要盖住文字 \\
交互状态 & 普通、hover、focus、checked 是否完整 & 只写 hover 会导致键盘 focus 不可见 \\
最终覆盖区 & 为什么必须放在文件末尾 & 它修复的是曲线图内部圆角，而不是重复布局 \\
\end{longtable}

示例应保持“一个代码块解决一个问题”。下面的代码只负责设置行的颜色和 focus，不应顺便改动页面宽度：

\begin{Shaded}
\begin{Verbatim}[breaklines=true,breakanywhere=true,fontsize=\small,baselinestretch=1.15]
/* 目标：设置行可读，键盘 focus 可见，且 hover 不改变布局尺寸。 */
html[data-theme="cute-blue"] [data-theme-ui="setting-row"] {
  background: var(--cute-blue-setting-row-background);
  border-color: var(--cute-blue-border);
}

html[data-theme="cute-blue"] [data-theme-ui="setting-row"]:focus-within {
  outline: 2px solid var(--cute-blue-focus-ring);
  outline-offset: 2px;
}
\end{Verbatim}
\end{Shaded}

不要把整份 \texttt{theme.css} 放进指南正文。完整主题应该放在模板目录，指南只保留能解释接口的最小片段；这样作者可以逐行理解，也不会在复制时带入与当前页面无关的规则。

\subsubsection{5.6 代码表格的排版和解释方法}

代码表格用于建立“代码入口 -> 视觉结果”的映射，而不是替代正文。表格中的代码列应保持左对齐、使用等宽字体；说明列使用短句，每格只表达一个动作。需要长篇解释时，应放到表格前后的段落中。

\begin{longtable}[]{@{}L{0.28\textwidth}L{0.17\textwidth}L{0.31\textwidth}L{0.18\textwidth}@{}}
\toprule\noalign{}\rowcolor{tablehead}
代码入口 & 所在层 & 主要作用 & 不应承担的工作 \\
\midrule\noalign{}
\endhead
\bottomrule\noalign{}
\endlastfoot
\texttt{html[data-theme="id"]} & 主题作用域 & 限制样式只作用于当前主题 & 不负责切换系统明暗 \\
\texttt{.glacier-shell} & 应用外壳 & 设置窗口背景和基础文字色 & 不重建页面 flex 布局 \\
\shortstack[l]{\texttt{data-theme-card=}\\\texttt{curve-editor}} & 曲线卡片 & 修正编辑器外层的视觉层级 & 不假定它就是内部绘图区 \\
\shortstack[l]{\texttt{data-theme-ui=}\\\texttt{setting-row}} & 交互组件 & 设置行背景、边框和 focus & 不在 hover 时修改高度 \\
\texttt{@font-face} & 资源层 & 注册授权的 WOFF2 子集 & 不加载远程字体 \\
\end{longtable}

表格行的顺序建议从稳定入口到细粒度入口：先写主题作用域，再写外壳类、页面/卡片钩子，最后写交互组件。这样读者可以从上到下建立覆盖层级，也能快速定位“规则没有生效”是作用域问题还是选择器问题。

\subsubsection{5.7 正文段落与对齐规则}

正文每段只讲一个判断或一个操作。推荐先给结论，再解释原因，最后补充限制；一段通常保持两到四句。不要把表格中的每一行再逐字重复成一串短段落，也不要用连续单行文本模拟表格。

代码、路径、属性值放在行内代码中，中文解释保持正常比例字体。路径很长时在目录分隔符处换行，不要为了“看起来一行”缩小字号。代码块使用固定的行间距，表格使用较宽的行高；两者都不应通过负边距压缩。

段落和表格之间留出明确的垂直间距。表格标题行使用轻微底色，正文行使用交替浅色或留白区分；颜色只用于层级提示，不能成为理解内容的唯一条件。

\begin{center}\rule{0.5\linewidth}{0.5pt}\end{center}

\subsection{\texorpdfstring{6. FanControl
外壳组件（\texttt{.glacier-*}）}{6. FanControl 外壳组件（.glacier-*）}}\label{6-fancontrol-ux5916ux58f3ux7ec4ux4ef6glacier-}

这些类是主题最稳定的入口。它们描述的是布局角色，而不是某一种颜色。

\begin{longtable}[]{@{}L{0.31\textwidth}L{0.25\textwidth}L{0.32\textwidth}@{}}
\toprule\noalign{}\rowcolor{tablehead}
类名 & 组件角色 & 常见用途 \\
\midrule\noalign{}
\endhead
\bottomrule\noalign{}
\endlastfoot
\themeclass{.glacier-shell} & 应用总外壳 & 背景、圆角、整体字体、窗口边界 \\
\themeclass{.glacier-titlebar} & 自定义标题栏 & 高度、透明度、底部描边 \\
\themeclass{.glacier-sidebar} & 左侧导航栏 & 背景、宽度适配、徽标、选中态 \\
\themeclass{.glacier-content} & 右侧内容外层 & 内容区布局和滚动外层 \\
\themeclass{.glacier-content-panel} & 实际滚动面板 &
页面背景、幕布、内容内边距 \\
\themeclass{.glacier-native-backdrop} & 原生材质存在时的外壳 &
透明度、系统背景兜底 \\
\themeclass{.glacier-hero-card} & 首页顶部大卡片 &
主视觉、设备名称、快捷操作 \\
\themeclass{.glacier-hero-content} & 大卡片文字和按钮区域 &
留出插画空间、控制层级 \\
\themeclass{.glacier-hero-actions} & 大卡片操作按钮区 &
对齐、胶囊容器、按钮间距 \\
\themeclass{.glacier-hero-art} & 首页插画容器 & 插画裁切、渐隐和装饰标签 \\
\themeclass{.glacier-hero-art-label} & 插画文字标签 &
小型装饰标题，不能遮挡主信息 \\
\themeclass{.glacier-operator-art} & 插画图片元素 &
\texttt{object-fit}、透明度和混合模式 \\
\themeclass{.glacier-metric-card} & 温度、转速等指标卡 & 背景、环形图、hover
阴影 \\
\themeclass{.glacier-control-card} & 控制与保护大卡片 &
模式、功耗、保护开关区域 \\
\themeclass{.glacier-stat-tile} & 控制卡内的小磁贴 &
当前模式、温度状态、功耗数值 \\
\themeclass{.glacier-chart-card} & 曲线和历史图卡片 &
图表外框、装饰、标题和状态 \\
\themeclass{.glacier-chart-canvas} & 图表绘制区域 &
网格、内阴影、扫描光和背景 \\
\themeclass{.glacier-page-card-fade} & 页面卡片进入动画 &
控制渐入，不要制造布局跳动 \\
\end{longtable}

\subsubsection{6.1 外壳示例}\label{61-ux5916ux58f3ux793aux4f8b}

\begin{Shaded}
\begin{Highlighting}
\CommentTok{/* 外壳只负责主题背景，不改变页面的 flex 和 overflow。 */}
\NormalTok{html}\ExtensionTok{[}\SpecialStringTok{data{-}theme}\OperatorTok{=}\StringTok{"cute{-}blue"}\ExtensionTok{]} \FunctionTok{.glacier{-}shell}\NormalTok{ \{}
  \KeywordTok{background}\CharTok{:} \FunctionTok{var(}\VariableTok{{-}{-}cute{-}blue{-}shell{-}background}\FunctionTok{)}\OperatorTok{;}
  \KeywordTok{color}\CharTok{:} \FunctionTok{var(}\VariableTok{{-}{-}foreground}\FunctionTok{)}\OperatorTok{;}
\NormalTok{\}}

\CommentTok{/* 侧栏是固定宽度区域，装饰不要改变它的尺寸。 */}
\NormalTok{html}\ExtensionTok{[}\SpecialStringTok{data{-}theme}\OperatorTok{=}\StringTok{"cute{-}blue"}\ExtensionTok{]} \FunctionTok{.glacier{-}sidebar}\NormalTok{ \{}
  \KeywordTok{background}\CharTok{:} \FunctionTok{var(}\VariableTok{{-}{-}cute{-}blue{-}sidebar{-}background}\FunctionTok{)}\OperatorTok{;}
  \KeywordTok{border{-}color}\CharTok{:} \FunctionTok{var(}\VariableTok{{-}{-}cute{-}blue{-}border}\FunctionTok{)}\OperatorTok{;}
\NormalTok{\}}

\CommentTok{/* 内容面板允许滚动，背景装饰应当使用伪元素而不是插入布局节点。 */}
\NormalTok{html}\ExtensionTok{[}\SpecialStringTok{data{-}theme}\OperatorTok{=}\StringTok{"cute{-}blue"}\ExtensionTok{]} \FunctionTok{.glacier{-}content{-}panel}\NormalTok{ \{}
  \KeywordTok{background}\CharTok{:} \FunctionTok{var(}\VariableTok{{-}{-}cute{-}blue{-}content{-}background}\FunctionTok{)}\OperatorTok{;}
\NormalTok{\}}
\end{Highlighting}
\end{Shaded}

解释：第一段改变总背景和文字颜色；第二段只改变侧栏的视觉，不改变
\texttt{width}；第三段只改变滚动区域的底色，避免背景装饰参与内容高度计算。

\begin{center}\rule{0.5\linewidth}{0.5pt}\end{center}

\subsection{\texorpdfstring{7.
页面级钩子：\texttt{data-theme-page}}{7. 页面级钩子：data-theme-page}}\label{7-ux9875ux9762ux7ea7ux94a9ux5b50data-theme-page}

FanControl 会在 \texttt{.glacier-shell}
上标记当前页面。页面钩子适合做大块背景和页面级视觉差异，不适合写通用按钮颜色。

当前常用页面值：

\begin{longtable}[]{@{}L{0.36\textwidth}L{0.58\textwidth}@{}}
\toprule\noalign{}\rowcolor{tablehead}
值 & 页面 \\
\midrule\noalign{}
\endhead
\bottomrule\noalign{}
\endlastfoot
\texttt{status} & 首页 / 设备状态 \\
\texttt{curve} & 风扇曲线和历史趋势 \\
\texttt{control} & 设置和控制面板 \\
\texttt{about} & 关于页 \\
\texttt{devices} & 高级设备信息页 \\
\end{longtable}

示例：

\begin{Shaded}
\begin{Highlighting}
\CommentTok{/* 首页可以有主视觉，但内容区仍然保持清晰。 */}
\NormalTok{html}\ExtensionTok{[}\SpecialStringTok{data{-}theme}\OperatorTok{=}\StringTok{"cute{-}blue"}\ExtensionTok{]} \FunctionTok{.glacier{-}shell}\ExtensionTok{[}\SpecialStringTok{data{-}theme{-}page}\OperatorTok{=}\StringTok{"status"}\ExtensionTok{]} \FunctionTok{.glacier{-}content{-}panel}\NormalTok{ \{}
  \KeywordTok{background}\CharTok{:} \FunctionTok{var(}\VariableTok{{-}{-}cute{-}blue{-}home{-}background}\FunctionTok{)}\OperatorTok{;}
\NormalTok{\}}

\CommentTok{/* 曲线页减少背景装饰，给图表留出稳定对比度。 */}
\NormalTok{html}\ExtensionTok{[}\SpecialStringTok{data{-}theme}\OperatorTok{=}\StringTok{"cute{-}blue"}\ExtensionTok{]} \FunctionTok{.glacier{-}shell}\ExtensionTok{[}\SpecialStringTok{data{-}theme{-}page}\OperatorTok{=}\StringTok{"curve"}\ExtensionTok{]} \FunctionTok{.glacier{-}content{-}panel}\NormalTok{ \{}
  \KeywordTok{background}\CharTok{:} \FunctionTok{var(}\VariableTok{{-}{-}cute{-}blue{-}curve{-}background}\FunctionTok{)}\OperatorTok{;}
\NormalTok{\}}
\end{Highlighting}
\end{Shaded}

旧版本如果没有
\texttt{data-theme-page}，这些规则会自然失效；因此页面钩子应该是增强层，不能是唯一的可读性来源。

\begin{center}\rule{0.5\linewidth}{0.5pt}\end{center}

\subsection{\texorpdfstring{8.
卡片级钩子：\texttt{data-theme-card}}{8. 卡片级钩子：data-theme-card}}\label{8-ux5361ux7247ux7ea7ux94a9ux5b50data-theme-card}

卡片钩子用于``同一页面中的一个语义卡片''。它比通用类更精确，但只在
FanControl 新版本中可靠。

\begin{longtable}[]{@{}L{0.36\textwidth}L{0.58\textwidth}@{}}
\toprule\noalign{}\rowcolor{tablehead}
值 & 组件 \\
\midrule\noalign{}
\endhead
\bottomrule\noalign{}
\endlastfoot
\texttt{settings-overview} & 设置页顶部设备概览 \\
\texttt{settings-overview-temperature} & 设置页顶部温度概览 \\
\texttt{settings-overview-device} & 设置页顶部设备概览 \\
\texttt{settings-offline-tip} & 设置页离线提示 \\
\texttt{device-hero} & 首页设备主信息卡 \\
\texttt{fan-curve-preview} & 首页小型曲线预览 \\
\texttt{temperature-history} & 首页温度 / 风扇历史图 \\
\texttt{cpu-temperature} & CPU 温度指标卡 \\
\texttt{gpu-temperature} & GPU 温度指标卡 \\
\texttt{fan-speed} & 风扇转速指标卡 \\
\texttt{control-mode} & 控制模式磁贴 \\
\texttt{temperature-state} & 温度状态磁贴 \\
\texttt{work-mode} & 工作模式磁贴 \\
\texttt{cpu-power} & CPU 功耗磁贴 \\
\texttt{gpu-power} & GPU 功耗磁贴 \\
\texttt{curve-header} & 曲线页标题和操作栏 \\
\texttt{curve-manual-gears} & 手动档位区域 \\
\texttt{curve-editor} & 曲线编辑器 \\
\texttt{curve-prediction} & 温度趋势预测区域 \\
\texttt{curve-learning} & 联合学习 / 学习控制区域 \\
\texttt{curve-history} & 历史趋势总卡片 \\
\texttt{curve-history-summary} & 历史摘要磁贴 \\
\texttt{curve-history-chart} & 历史图表卡片 \\
\end{longtable}

示例：

\begin{Shaded}
\begin{Highlighting}
\CommentTok{/* 只让 CPU 和 GPU 功耗磁贴使用强调色，避免影响其他统计项。 */}
\NormalTok{html}\ExtensionTok{[}\SpecialStringTok{data{-}theme}\OperatorTok{=}\StringTok{"cute{-}blue"}\ExtensionTok{]} \ExtensionTok{[}\SpecialStringTok{data{-}theme{-}card}\OperatorTok{=}\StringTok{"cpu{-}power"}\ExtensionTok{]}\OperatorTok{,}
\NormalTok{html}\ExtensionTok{[}\SpecialStringTok{data{-}theme}\OperatorTok{=}\StringTok{"cute{-}blue"}\ExtensionTok{]} \ExtensionTok{[}\SpecialStringTok{data{-}theme{-}card}\OperatorTok{=}\StringTok{"gpu{-}power"}\ExtensionTok{]}\NormalTok{ \{}
  \KeywordTok{border{-}color}\CharTok{:} \FunctionTok{var(}\VariableTok{{-}{-}cute{-}blue{-}power{-}border}\FunctionTok{)}\OperatorTok{;}
  \KeywordTok{background}\CharTok{:} \FunctionTok{var(}\VariableTok{{-}{-}cute{-}blue{-}power{-}background}\FunctionTok{)}\OperatorTok{;}
\NormalTok{\}}

\CommentTok{/* 曲线编辑器外层有时只是布局锚点，真正的圆角卡片在内部。 */}
\NormalTok{html}\ExtensionTok{[}\SpecialStringTok{data{-}theme}\OperatorTok{=}\StringTok{"cute{-}blue"}\ExtensionTok{]} \ExtensionTok{[}\SpecialStringTok{data{-}theme{-}card}\OperatorTok{=}\StringTok{"curve{-}editor"}\ExtensionTok{]} \OperatorTok{\textgreater{}} \FunctionTok{.relative.rounded{-}3xl}\NormalTok{ \{}
  \KeywordTok{border{-}radius}\CharTok{:} \DecValTok{1.25}\DataTypeTok{rem}\OperatorTok{;}
  \KeywordTok{overflow}\CharTok{:} \DecValTok{hidden}\OperatorTok{;}
\NormalTok{\}}
\end{Highlighting}
\end{Shaded}

第二条规则解释了一个常见问题：如果只给外层设置背景，内部图表卡片可能仍然出现方角；如果只给内部设置圆角，外层装饰又可能露出原生窗口背景。需要同时检查外层和内部真实绘图区。

\begin{center}\rule{0.5\linewidth}{0.5pt}\end{center}

\subsection{\texorpdfstring{9.
组件级钩子：\texttt{data-theme-ui}}{9. 组件级钩子：data-theme-ui}}\label{9-ux7ec4ux4ef6ux7ea7ux94a9ux5b50data-theme-ui}

UI
钩子用于按钮、输入、设置行、切换器和弹窗等交互组件。它们适合调节边框、hover、focus、选中态和可读性。

\begin{longtable}[]{@{}L{0.36\textwidth}L{0.58\textwidth}@{}}
\toprule\noalign{}\rowcolor{tablehead}
值 & 组件 \\
\midrule\noalign{}
\endhead
\bottomrule\noalign{}
\endlastfoot
\texttt{brand-mark} & 品牌或左上角标记 \\
\texttt{sidebar-item} & 左侧导航按钮 \\
\texttt{settings-tabs} & 设置页分区标签栏 \\
\texttt{settings-tab} & 设置页单个分区标签 \\
\texttt{settings-panels} & 设置页分区内容容器 \\
\texttt{settings-panel} & 设置页单个内容面板 \\
\texttt{setting-section} & 设置分区外框 \\
\texttt{setting-section-header} & 设置分区标题栏 \\
\texttt{setting-row} & 设置行 \\
\texttt{setting-row-icon} & 设置行左侧图标 \\
\texttt{setting-row-control} & 设置行右侧控制区 \\
\texttt{connection-panel} & 设备连接面板 \\
\texttt{compatibility-submenu} & 兼容模式子菜单 \\
\texttt{compatibility-submenu-row} & 兼容模式子菜单行 \\
\texttt{compatibility-nested} & 兼容模式嵌套内容 \\
\texttt{compatibility-nested-row} & 嵌套设置行 \\
\texttt{compatibility-nested-panel} & 嵌套控制面板 \\
\texttt{select-trigger} & 下拉选择器按钮 \\
\texttt{switch} & 开关轨道 \\
\texttt{switch-thumb} & 开关滑块 \\
\texttt{slider} & 滑块容器 \\
\texttt{manual-gear-dot} & 手动档位视觉节点 \\
\texttt{learning-target-temp} & 学习目标温度控制 \\
\texttt{curve-hints} & 曲线页提示标签 \\
\texttt{history-display-dialog} & 历史显示弹窗 \\
\texttt{history-display-row} & 历史显示项 \\
\texttt{hero-actions} & 首页主卡片操作区 \\
\end{longtable}

\subsubsection{9.1 设置行示例}\label{91-ux8bbeux7f6eux884cux793aux4f8b}

\begin{Shaded}
\begin{Highlighting}
\CommentTok{/* 设置行默认保持清楚的文字和边界。 */}
\NormalTok{html}\ExtensionTok{[}\SpecialStringTok{data{-}theme}\OperatorTok{=}\StringTok{"cute{-}blue"}\ExtensionTok{]} \ExtensionTok{[}\SpecialStringTok{data{-}theme{-}ui}\OperatorTok{=}\StringTok{"setting{-}row"}\ExtensionTok{]}\NormalTok{ \{}
  \KeywordTok{background}\CharTok{:} \FunctionTok{var(}\VariableTok{{-}{-}cute{-}blue{-}setting{-}row{-}background}\FunctionTok{)}\OperatorTok{;}
  \KeywordTok{border{-}color}\CharTok{:} \FunctionTok{var(}\VariableTok{{-}{-}cute{-}blue{-}border}\FunctionTok{)}\OperatorTok{;}
\NormalTok{\}}

\CommentTok{/* hover 只增加层次，不移动行，防止设置页滚动位置跳动。 */}
\NormalTok{html}\ExtensionTok{[}\SpecialStringTok{data{-}theme}\OperatorTok{=}\StringTok{"cute{-}blue"}\ExtensionTok{]} \ExtensionTok{[}\SpecialStringTok{data{-}theme{-}ui}\OperatorTok{=}\StringTok{"setting{-}row"}\ExtensionTok{]}\InformationTok{:hover}\NormalTok{ \{}
  \KeywordTok{background}\CharTok{:} \FunctionTok{var(}\VariableTok{{-}{-}cute{-}blue{-}setting{-}row{-}hover}\FunctionTok{)}\OperatorTok{;}
  \KeywordTok{box{-}shadow}\CharTok{:} \DecValTok{inset} \DecValTok{2}\DataTypeTok{px} \DecValTok{0} \DecValTok{0} \FunctionTok{var(}\VariableTok{{-}{-}cute{-}blue{-}accent}\FunctionTok{)}\OperatorTok{;}
\NormalTok{\}}
\end{Highlighting}
\end{Shaded}

不要在 hover 时修改 \texttt{padding}、\texttt{margin}、\texttt{height}
或 \texttt{border-width}，这些变化会造成整页布局抖动。

\subsubsection{9.2
下拉框和开关示例}\label{92-ux4e0bux62c9ux6846ux548cux5f00ux5173ux793aux4f8b}

\begin{Shaded}
\begin{Highlighting}
\CommentTok{/* 下拉触发按钮要同时处理普通、hover、打开和 focus{-}visible 状态。 */}
\NormalTok{html}\ExtensionTok{[}\SpecialStringTok{data{-}theme}\OperatorTok{=}\StringTok{"cute{-}blue"}\ExtensionTok{]} \ExtensionTok{[}\SpecialStringTok{data{-}theme{-}ui}\OperatorTok{=}\StringTok{"select{-}trigger"}\ExtensionTok{]}\NormalTok{ \{}
  \KeywordTok{background}\CharTok{:} \FunctionTok{var(}\VariableTok{{-}{-}cute{-}blue{-}control{-}background}\FunctionTok{)}\OperatorTok{;}
  \KeywordTok{border{-}color}\CharTok{:} \FunctionTok{var(}\VariableTok{{-}{-}cute{-}blue{-}control{-}border}\FunctionTok{)}\OperatorTok{;}
  \KeywordTok{color}\CharTok{:} \FunctionTok{var(}\VariableTok{{-}{-}foreground}\FunctionTok{)}\OperatorTok{;}
\NormalTok{\}}

\NormalTok{html}\ExtensionTok{[}\SpecialStringTok{data{-}theme}\OperatorTok{=}\StringTok{"cute{-}blue"}\ExtensionTok{]} \ExtensionTok{[}\SpecialStringTok{data{-}theme{-}ui}\OperatorTok{=}\StringTok{"select{-}trigger"}\ExtensionTok{]}\InformationTok{:hover}\OperatorTok{,}
\NormalTok{html}\ExtensionTok{[}\SpecialStringTok{data{-}theme}\OperatorTok{=}\StringTok{"cute{-}blue"}\ExtensionTok{]} \ExtensionTok{[}\SpecialStringTok{data{-}theme{-}ui}\OperatorTok{=}\StringTok{"select{-}trigger"}\ExtensionTok{][}\SpecialStringTok{data{-}state}\OperatorTok{=}\StringTok{"open"}\ExtensionTok{]}\NormalTok{ \{}
  \KeywordTok{border{-}color}\CharTok{:} \FunctionTok{var(}\VariableTok{{-}{-}cute{-}blue{-}accent}\FunctionTok{)}\OperatorTok{;}
\NormalTok{\}}

\NormalTok{html}\ExtensionTok{[}\SpecialStringTok{data{-}theme}\OperatorTok{=}\StringTok{"cute{-}blue"}\ExtensionTok{]} \ExtensionTok{[}\SpecialStringTok{data{-}theme{-}ui}\OperatorTok{=}\StringTok{"select{-}trigger"}\ExtensionTok{]}\InformationTok{:focus{-}visible}\NormalTok{ \{}
  \KeywordTok{outline}\CharTok{:} \DecValTok{2}\DataTypeTok{px} \DecValTok{solid} \FunctionTok{var(}\VariableTok{{-}{-}cute{-}blue{-}focus{-}ring}\FunctionTok{)}\OperatorTok{;}
  \KeywordTok{outline{-}offset}\CharTok{:} \DecValTok{2}\DataTypeTok{px}\OperatorTok{;}
\NormalTok{\}}

\CommentTok{/* 开关要同时覆盖未选中、选中和滑块颜色。 */}
\NormalTok{html}\ExtensionTok{[}\SpecialStringTok{data{-}theme}\OperatorTok{=}\StringTok{"cute{-}blue"}\ExtensionTok{]} \ExtensionTok{[}\SpecialStringTok{data{-}theme{-}ui}\OperatorTok{=}\StringTok{"switch"}\ExtensionTok{]}\NormalTok{ \{}
  \KeywordTok{background}\CharTok{:} \FunctionTok{var(}\VariableTok{{-}{-}cute{-}blue{-}switch{-}off}\FunctionTok{)}\OperatorTok{;}
\NormalTok{\}}

\NormalTok{html}\ExtensionTok{[}\SpecialStringTok{data{-}theme}\OperatorTok{=}\StringTok{"cute{-}blue"}\ExtensionTok{]} \ExtensionTok{[}\SpecialStringTok{data{-}theme{-}ui}\OperatorTok{=}\StringTok{"switch"}\ExtensionTok{][}\SpecialStringTok{data{-}state}\OperatorTok{=}\StringTok{"checked"}\ExtensionTok{]}\NormalTok{ \{}
  \KeywordTok{background}\CharTok{:} \FunctionTok{var(}\VariableTok{{-}{-}cute{-}blue{-}switch{-}on}\FunctionTok{)}\OperatorTok{;}
\NormalTok{\}}

\NormalTok{html}\ExtensionTok{[}\SpecialStringTok{data{-}theme}\OperatorTok{=}\StringTok{"cute{-}blue"}\ExtensionTok{]} \ExtensionTok{[}\SpecialStringTok{data{-}theme{-}ui}\OperatorTok{=}\StringTok{"switch{-}thumb"}\ExtensionTok{]}\NormalTok{ \{}
  \KeywordTok{background}\CharTok{:} \FunctionTok{var(}\VariableTok{{-}{-}cute{-}blue{-}switch{-}thumb}\FunctionTok{)}\OperatorTok{;}
\NormalTok{\}}
\end{Highlighting}
\end{Shaded}

\begin{center}\rule{0.5\linewidth}{0.5pt}\end{center}

\subsection{10. Basic
模板变量逐组说明}\label{10-basic-ux6a21ux677fux53d8ux91cfux9010ux7ec4ux8bf4ux660e}

Basic
模板的变量区应该是主题最先修改的地方。变量名以主题模板中的实际前缀为准；复制后建议将
\texttt{-\/-basic-} 改为自己的前缀，例如 \texttt{-\/-cute-blue-}。

\subsubsection{10.1 变量分类总览}

主题变量不是同一种东西。控制颜色决定文字、边框、按钮和状态是否可读；纹理变量只负责材质、渐变和图片；几何变量决定尺寸和层级。先判断变量类型，再决定它应该放在控制组件还是装饰层。

\begin{longtable}[]{@{}L{0.15\textwidth}L{0.34\textwidth}L{0.22\textwidth}L{0.22\textwidth}@{}}
\toprule\noalign{}\rowcolor{tablehead}
变量类别 & 典型变量 & CSS 值类型 & 负责什么和可感知范围 \\
\midrule\noalign{}
\endhead
\bottomrule\noalign{}
\endlastfoot
控制颜色 & \texttt{-\/-background}、\texttt{-\/-foreground}、\texttt{-\/-card}、\texttt{-\/-border}、\texttt{-\/-primary} & \texttt{<color>} & 页面、文字、卡片、边框和主操作按钮；影响所有页面和交互控件。 \\
状态颜色 & \texttt{-\/-primary-foreground}、\texttt{-\/-muted-foreground}、\texttt{-\/-focus-ring}、\texttt{-\/-switch-on} & \texttt{<color>} 或带透明度的 \texttt{<color>} & hover、focus、checked、禁用和错误状态；影响键盘操作、开关、下拉框和提示。 \\
图表颜色 & \texttt{-\/-chart-1} 至 \texttt{-\/-chart-4} & \texttt{<color>} & CPU、GPU、风扇和功耗趋势；影响曲线页、历史图和悬浮提示。 \\
纹理/材质 & \texttt{-\/-basic-sidebar-background}、\texttt{-\/-advanced-hero-image}、\texttt{-\/-advanced-curtain-image} & \texttt{<image>}、\texttt{linear-gradient()} 或 \texttt{background} 简写 & 渐变、插画、幕布和背景质感；只影响首页、侧栏和内容面板，不改变控件语义。 \\
透明保护层 & \texttt{-\/-advanced-hero-overlay}、\texttt{-\/-dialog-overlay} & 带 alpha 的 \texttt{<color>} & 在图片上方保护文字和按钮对比度；影响插画、弹窗和下拉菜单。 \\
阴影层级 & \texttt{-\/-basic-dialog-shadow}、\texttt{-\/-card-shadow} & \texttt{<shadow>} & 表达卡片、弹窗和浮层的前后关系；影响边界、弹窗和 hover。 \\
几何/层级 & \texttt{-\/-radius}、\texttt{-\/-radius-sm}、\texttt{-\/-z-content}、\texttt{-\/-z-overlay} & \texttt{<length>} 或 \texttt{<integer>} & 控制圆角、层级和装饰覆盖顺序；影响所有页面的布局稳定性。 \\
资源/文字 & \texttt{-\/-basic-sidebar-logo-image}、\texttt{-\/-basic-sidebar-logo-text} & \texttt{<image>} 或 \texttt{<string>} & 注册侧栏图片徽标和文字徽标；影响左侧导航和品牌标记。 \\
\end{longtable}

变量后缀也有固定含义。\texttt{-\/-background} 通常是表面颜色或材质，\texttt{-\/-foreground} 是表面上的文字/图标，\texttt{-\/-border} 是边界，\texttt{-\/-hover}、\texttt{-\/-active}、\texttt{-\/-checked} 和 \texttt{-\/-focus-ring} 是状态，\texttt{-\/-image} 是本地资源，\texttt{-\/-overlay} 是保护层，\texttt{-\/-shadow} 是阴影，\texttt{-\/-radius} 和 \texttt{-\/-z-*} 是结构参数。不要把图片写进 \texttt{-\/-primary}，也不要把按钮文字色写进 \texttt{-\/-hero-image}。

控制颜色变量必须始终有可读的配对关系。例如 \texttt{-\/-primary} 负责按钮背景，\texttt{-\/-primary-foreground} 负责按钮文字；\texttt{-\/-card} 负责卡片表面，\texttt{-\/-card-foreground} 负责卡片文字；\texttt{-\/-background} 和 \texttt{-\/-foreground} 负责页面默认层。纹理变量可以被删除或替换，但删除后必须退回到纯色背景。

\begin{Shaded}
\begin{Verbatim}[breaklines=true,breakanywhere=true,fontsize=\small,baselinestretch=1.15]
html[data-theme="cute-blue"] {
  /* 控制颜色：直接影响文字、按钮和边框。 */
  --primary: #2f6df6;
  --primary-foreground: #ffffff;
  --border: #d8e1ef;

  /* 纹理/材质：只影响背景，不改变控件语义。 */
  --hero-image: url("hero.webp");
  --hero-overlay: rgb(255 255 255 / 65%);
}
\end{Verbatim}
\end{Shaded}

\subsubsection{10.2 基础颜色}\label{102-ux57faux7840ux989cux8272}

\begin{Shaded}
\begin{Highlighting}
\NormalTok{html}\ExtensionTok{[}\SpecialStringTok{data{-}theme}\OperatorTok{=}\StringTok{"basic"}\ExtensionTok{]}\NormalTok{ \{}
  \VariableTok{{-}{-}background}\CharTok{:} \ConstantTok{\#f6f8fc}\OperatorTok{;}       \CommentTok{/* 页面底色。 */}
  \VariableTok{{-}{-}foreground}\CharTok{:} \ConstantTok{\#172033}\OperatorTok{;}       \CommentTok{/* 主文字颜色。 */}
  \VariableTok{{-}{-}card}\CharTok{:} \ConstantTok{\#ffffff}\OperatorTok{;}             \CommentTok{/* 卡片底色。 */}
  \VariableTok{{-}{-}card{-}foreground}\CharTok{:} \ConstantTok{\#172033}\OperatorTok{;}  \CommentTok{/* 卡片内文字。 */}
  \VariableTok{{-}{-}muted}\CharTok{:} \ConstantTok{\#eaf0f8}\OperatorTok{;}            \CommentTok{/* 次级背景和禁用区域。 */}
  \VariableTok{{-}{-}muted{-}foreground}\CharTok{:} \ConstantTok{\#65738a}\OperatorTok{;} \CommentTok{/* 次级文字。 */}
  \VariableTok{{-}{-}border}\CharTok{:} \ConstantTok{\#d8e1ef}\OperatorTok{;}           \CommentTok{/* 通用边框。 */}
  \VariableTok{{-}{-}primary}\CharTok{:} \ConstantTok{\#2f6df6}\OperatorTok{;}          \CommentTok{/* 主操作色。 */}
  \VariableTok{{-}{-}primary{-}foreground}\CharTok{:} \ConstantTok{\#ffffff}\OperatorTok{;} \CommentTok{/* 主操作按钮文字。 */}
\NormalTok{\}}
\end{Highlighting}
\end{Shaded}

不要只改变 \texttt{-\/-primary} 而不检查
\texttt{-\/-primary-foreground}。浅色按钮配浅色文字或深色按钮配深色文字都会产生低对比度。

\subsubsection{10.3 图表颜色}\label{103-ux56feux8868ux989cux8272}

功耗、CPU 温度、GPU 温度和风扇趋势可能同时出现。颜色必须彼此区分：

\begin{Shaded}
\begin{Highlighting}
\NormalTok{html}\ExtensionTok{[}\SpecialStringTok{data{-}theme}\OperatorTok{=}\StringTok{"basic"}\ExtensionTok{]}\NormalTok{ \{}
  \VariableTok{{-}{-}chart{-}1}\CharTok{:} \ConstantTok{\#2f6df6}\OperatorTok{;} \CommentTok{/* CPU 或主曲线。 */}
  \VariableTok{{-}{-}chart{-}2}\CharTok{:} \ConstantTok{\#e05252}\OperatorTok{;} \CommentTok{/* GPU 或第二条温度曲线。 */}
  \VariableTok{{-}{-}chart{-}3}\CharTok{:} \ConstantTok{\#2b9b67}\OperatorTok{;} \CommentTok{/* 风扇或状态曲线。 */}
  \VariableTok{{-}{-}chart{-}4}\CharTok{:} \ConstantTok{\#d59025}\OperatorTok{;} \CommentTok{/* 功耗或辅助曲线。 */}
\NormalTok{\}}
\end{Highlighting}
\end{Shaded}

不要把所有曲线都设置成同一色相的浅色版本。相邻曲线在悬浮提示和截图中必须能快速区分。

\subsubsection{10.4
圆角和层级}\label{104-ux5706ux89d2ux548cux5c42ux7ea7}

\begin{Shaded}
\begin{Highlighting}
\NormalTok{html}\ExtensionTok{[}\SpecialStringTok{data{-}theme}\OperatorTok{=}\StringTok{"basic"}\ExtensionTok{]}\NormalTok{ \{}
  \VariableTok{{-}{-}radius}\CharTok{:} \DecValTok{0.75}\DataTypeTok{rem}\OperatorTok{;}       \CommentTok{/* 通用圆角。 */}
  \VariableTok{{-}{-}radius{-}sm}\CharTok{:} \DecValTok{0.5}\DataTypeTok{rem}\OperatorTok{;}     \CommentTok{/* 小控件圆角。 */}
  \VariableTok{{-}{-}z{-}content}\CharTok{:} \DecValTok{1}\OperatorTok{;}           \CommentTok{/* 内容层。 */}
  \VariableTok{{-}{-}z{-}decoration}\CharTok{:} \DecValTok{2}\OperatorTok{;}        \CommentTok{/* 装饰层。 */}
  \VariableTok{{-}{-}z{-}overlay}\CharTok{:} \DecValTok{20}\OperatorTok{;}          \CommentTok{/* 弹窗和浮层。 */}
\NormalTok{\}}
\end{Highlighting}
\end{Shaded}

装饰层不应该高于弹窗、下拉框和
toast。层级变量是基础模板中的安全栏，不建议随意增大。

\subsubsection{10.5
侧栏、内容区和弹窗}\label{105-ux4fa7ux680fux5185ux5bb9ux533aux548cux5f39ux7a97}

\begin{Shaded}
\begin{Highlighting}
\NormalTok{html}\ExtensionTok{[}\SpecialStringTok{data{-}theme}\OperatorTok{=}\StringTok{"basic"}\ExtensionTok{]}\NormalTok{ \{}
  \VariableTok{{-}{-}basic{-}sidebar{-}background}\CharTok{:} \FunctionTok{linear{-}gradient(}\DecValTok{180}\DataTypeTok{deg}\OperatorTok{,} \ConstantTok{\#ffffff}\OperatorTok{,} \ConstantTok{\#eef5ff}\FunctionTok{)}\OperatorTok{;}
  \VariableTok{{-}{-}basic{-}sidebar{-}foreground}\CharTok{:} \ConstantTok{\#26364f}\OperatorTok{;}
  \VariableTok{{-}{-}basic{-}content{-}background}\CharTok{:} \ConstantTok{\#f6f8fc}\OperatorTok{;}
  \VariableTok{{-}{-}basic{-}dialog{-}background}\CharTok{:} \ConstantTok{\#ffffff}\OperatorTok{;}
  \VariableTok{{-}{-}basic{-}dialog{-}border}\CharTok{:} \ConstantTok{\#c9d5e7}\OperatorTok{;}
  \VariableTok{{-}{-}basic{-}dialog{-}shadow}\CharTok{:} \DecValTok{0} \DecValTok{18}\DataTypeTok{px} \DecValTok{60}\DataTypeTok{px} \FunctionTok{rgb(}\DecValTok{29} \DecValTok{52} \DecValTok{87}\NormalTok{ / }\DecValTok{18}\DataTypeTok{\%}\FunctionTok{)}\OperatorTok{;}
\NormalTok{\}}
\end{Highlighting}
\end{Shaded}

背景可以使用渐变，但正文卡片应该保持足够不透明。弹窗和下拉菜单尤其不能透到难以读取。

\subsubsection{10.6 徽标变量}\label{106-ux5fbdux6807ux53d8ux91cf}

\begin{Shaded}
\begin{Highlighting}
\NormalTok{html}\ExtensionTok{[}\SpecialStringTok{data{-}theme}\OperatorTok{=}\StringTok{"basic"}\ExtensionTok{]}\NormalTok{ \{}
  \VariableTok{{-}{-}basic{-}sidebar{-}logo{-}text}\CharTok{:} \StringTok{"Basic"}\OperatorTok{;} \CommentTok{/* 文字徽标。 */}
  \VariableTok{{-}{-}basic{-}sidebar{-}logo{-}image}\CharTok{:} \DecValTok{none}\OperatorTok{;}    \CommentTok{/* 图片徽标；有图片时写 url。 */}
\NormalTok{\}}
\end{Highlighting}
\end{Shaded}

二选一最稳定。若同时保留文字和图片，必须确认二者不会重叠。

\begin{center}\rule{0.5\linewidth}{0.5pt}\end{center}

\subsection{11. Advanced
模板的页面和资源设计}\label{11-advanced-ux6a21ux677fux7684ux9875ux9762ux548cux8d44ux6e90ux8bbeux8ba1}

Advanced
不是``把所有装饰都打开''。它仍然需要先建立一个可读的结构，再逐层增加视觉。

\subsubsection{11.1
首页主视觉}\label{111-ux9996ux9875ux4e3bux89c6ux89c9}

\begin{Shaded}
\begin{Highlighting}
\NormalTok{html}\ExtensionTok{[}\SpecialStringTok{data{-}theme}\OperatorTok{=}\StringTok{"advanced"}\ExtensionTok{]}\NormalTok{ \{}
  \VariableTok{{-}{-}advanced{-}hero{-}image}\CharTok{:} \FunctionTok{url(}\StringTok{"hero.webp"}\FunctionTok{)}\OperatorTok{;} \CommentTok{/* 首页主插画。 */}
  \VariableTok{{-}{-}advanced{-}hero{-}overlay}\CharTok{:} \FunctionTok{rgb(}\DecValTok{255} \DecValTok{255} \DecValTok{255}\NormalTok{ / }\DecValTok{65}\DataTypeTok{\%}\FunctionTok{)}\OperatorTok{;} \CommentTok{/* 文字保护层。 */}
\NormalTok{\}}

\NormalTok{html}\ExtensionTok{[}\SpecialStringTok{data{-}theme}\OperatorTok{=}\StringTok{"advanced"}\ExtensionTok{]} \FunctionTok{.glacier{-}hero{-}art}\InformationTok{::before}\NormalTok{ \{}
  \KeywordTok{content}\CharTok{:} \StringTok{""}\OperatorTok{;} \CommentTok{/* 伪元素必须有内容，才能绘制背景。 */}
  \KeywordTok{position}\CharTok{:} \DecValTok{absolute}\OperatorTok{;} \CommentTok{/* 让插画脱离文字布局。 */}
  \KeywordTok{inset}\CharTok{:} \DecValTok{0}\OperatorTok{;} \CommentTok{/* 填满插画容器。 */}
  \KeywordTok{background}\CharTok{:} \FunctionTok{var(}\VariableTok{{-}{-}advanced{-}hero{-}image}\FunctionTok{)} \DecValTok{right} \DecValTok{center}\NormalTok{ / }\DecValTok{contain} \DecValTok{no{-}repeat}\OperatorTok{;}
  \KeywordTok{opacity}\CharTok{:} \DecValTok{0.72}\OperatorTok{;} \CommentTok{/* 降低插画对文字的干扰。 */}
  \KeywordTok{pointer{-}events}\CharTok{:} \DecValTok{none}\OperatorTok{;} \CommentTok{/* 装饰不能拦截按钮点击。 */}
\NormalTok{\}}
\end{Highlighting}
\end{Shaded}

首页插画应该放在 \texttt{glacier-hero-art}
中，并给文字区域保留足够空间。不要把高对比人物脸放在设备名称或操作按钮正后方。

\subsubsection{11.2 背景幕布}\label{112-ux80ccux666fux5e55ux5e03}

\begin{Shaded}
\begin{Highlighting}
\NormalTok{html}\ExtensionTok{[}\SpecialStringTok{data{-}theme}\OperatorTok{=}\StringTok{"advanced"}\ExtensionTok{]}\NormalTok{ \{}
  \VariableTok{{-}{-}advanced{-}curtain{-}image}\CharTok{:} \FunctionTok{url(}\StringTok{"assets/stickers/curtain.webp"}\FunctionTok{)}\OperatorTok{;}
\NormalTok{\}}

\NormalTok{html}\ExtensionTok{[}\SpecialStringTok{data{-}theme}\OperatorTok{=}\StringTok{"advanced"}\ExtensionTok{]} \FunctionTok{.glacier{-}content{-}panel}\InformationTok{::before}\NormalTok{ \{}
  \KeywordTok{content}\CharTok{:} \StringTok{""}\OperatorTok{;} \CommentTok{/* 创建幕布层。 */}
  \KeywordTok{position}\CharTok{:} \DecValTok{absolute}\OperatorTok{;} \CommentTok{/* 不参与内容高度。 */}
  \KeywordTok{inset}\CharTok{:} \DecValTok{0}\OperatorTok{;} \CommentTok{/* 覆盖整个内容面板。 */}
  \KeywordTok{background}\CharTok{:} \FunctionTok{var(}\VariableTok{{-}{-}advanced{-}curtain{-}image}\FunctionTok{)} \DecValTok{top} \DecValTok{center}\NormalTok{ / }\DecValTok{cover} \DecValTok{no{-}repeat}\OperatorTok{;}
  \KeywordTok{opacity}\CharTok{:} \DecValTok{0.18}\OperatorTok{;} \CommentTok{/* 幕布只能提供质感。 */}
  \KeywordTok{pointer{-}events}\CharTok{:} \DecValTok{none}\OperatorTok{;} \CommentTok{/* 不挡滚动和点击。 */}
\NormalTok{\}}
\end{Highlighting}
\end{Shaded}

内容面板需要有 \texttt{position:\ relative} 和合适的
z-index，正文需要位于伪元素上方。背景幕布在设置页通常应该比首页更淡。

\subsubsection{11.3
贴纸和小装饰}\label{113-ux8d34ux7eb8ux548cux5c0fux88c5ux9970}

\begin{Shaded}
\begin{Highlighting}
\NormalTok{html}\ExtensionTok{[}\SpecialStringTok{data{-}theme}\OperatorTok{=}\StringTok{"advanced"}\ExtensionTok{]}\NormalTok{ \{}
  \VariableTok{{-}{-}advanced{-}star}\CharTok{:} \FunctionTok{url(}\StringTok{"decorations/star.svg"}\FunctionTok{)}\OperatorTok{;}
\NormalTok{\}}

\NormalTok{html}\ExtensionTok{[}\SpecialStringTok{data{-}theme}\OperatorTok{=}\StringTok{"advanced"}\ExtensionTok{]} \FunctionTok{.glacier{-}metric{-}card}\InformationTok{::after}\NormalTok{ \{}
  \KeywordTok{content}\CharTok{:} \StringTok{""}\OperatorTok{;} \CommentTok{/* 创建装饰层。 */}
  \KeywordTok{position}\CharTok{:} \DecValTok{absolute}\OperatorTok{;} \CommentTok{/* 不改变卡片尺寸。 */}
  \KeywordTok{right}\CharTok{:} \DecValTok{0.75}\DataTypeTok{rem}\OperatorTok{;} \CommentTok{/* 与卡片右边保持距离。 */}
  \KeywordTok{bottom}\CharTok{:} \DecValTok{0.75}\DataTypeTok{rem}\OperatorTok{;} \CommentTok{/* 与卡片底部保持距离。 */}
  \KeywordTok{width}\CharTok{:} \DecValTok{2}\DataTypeTok{rem}\OperatorTok{;} \CommentTok{/* 固定尺寸，防止布局抖动。 */}
  \KeywordTok{height}\CharTok{:} \DecValTok{2}\DataTypeTok{rem}\OperatorTok{;}
  \KeywordTok{background}\CharTok{:} \FunctionTok{var(}\VariableTok{{-}{-}advanced{-}star}\FunctionTok{)} \DecValTok{center}\NormalTok{ / }\DecValTok{contain} \DecValTok{no{-}repeat}\OperatorTok{;}
  \KeywordTok{opacity}\CharTok{:} \DecValTok{0.35}\OperatorTok{;} \CommentTok{/* 装饰不应盖过数值。 */}
  \KeywordTok{pointer{-}events}\CharTok{:} \DecValTok{none}\OperatorTok{;}
\NormalTok{\}}
\end{Highlighting}
\end{Shaded}

伪元素装饰必须有固定尺寸、低透明度和
\texttt{pointer-events:\ none}。它不能影响卡片的实际内容高度。

\subsubsection{11.4
卡片透明度}\label{114-ux5361ux7247ux900fux660eux5ea6}

\begin{Shaded}
\begin{Highlighting}
\NormalTok{html}\ExtensionTok{[}\SpecialStringTok{data{-}theme}\OperatorTok{=}\StringTok{"advanced"}\ExtensionTok{]} \ExtensionTok{[}\SpecialStringTok{data{-}theme{-}card}\OperatorTok{=}\StringTok{"temperature{-}history"}\ExtensionTok{]}\NormalTok{ \{}
  \KeywordTok{background}\CharTok{:} \FunctionTok{rgb(}\DecValTok{255} \DecValTok{255} \DecValTok{255}\NormalTok{ / }\DecValTok{82}\DataTypeTok{\%}\FunctionTok{)}\OperatorTok{;} \CommentTok{/* 图表背景保持高可读性。 */}
  \KeywordTok{border{-}color}\CharTok{:} \FunctionTok{rgb(}\DecValTok{76} \DecValTok{105} \DecValTok{145}\NormalTok{ / }\DecValTok{25}\DataTypeTok{\%}\FunctionTok{)}\OperatorTok{;} \CommentTok{/* 边框只保留层次。 */}
  \KeywordTok{backdrop{-}filter}\CharTok{:} \FunctionTok{blur(}\DecValTok{14}\DataTypeTok{px}\FunctionTok{)}\OperatorTok{;} \CommentTok{/* 背景复杂时只使用中等 blur。 */}
\NormalTok{\}}
\end{Highlighting}
\end{Shaded}

透明度低于约 60\%
时，图表网格和文字容易与背景混在一起。设置页、曲线编辑器和历史图通常应使用更高不透明度。

\begin{center}\rule{0.5\linewidth}{0.5pt}\end{center}

\subsection{12.
左侧徽标和品牌标记}\label{12-ux5de6ux4fa7ux5fbdux6807ux548cux54c1ux724cux6807ux8bb0}

\subsubsection{12.1 文字徽标}\label{121-ux6587ux5b57ux5fbdux6807}

\begin{Shaded}
\begin{Highlighting}
\NormalTok{html}\ExtensionTok{[}\SpecialStringTok{data{-}theme}\OperatorTok{=}\StringTok{"basic"}\ExtensionTok{]}\NormalTok{ \{}
  \VariableTok{{-}{-}basic{-}sidebar{-}logo{-}text}\CharTok{:} \StringTok{"AUX"}\OperatorTok{;} \CommentTok{/* 文字内容。 */}
\NormalTok{\}}

\NormalTok{html}\ExtensionTok{[}\SpecialStringTok{data{-}theme}\OperatorTok{=}\StringTok{"basic"}\ExtensionTok{]} \FunctionTok{.glacier{-}sidebar}\InformationTok{::after}\NormalTok{ \{}
  \KeywordTok{content}\CharTok{:} \FunctionTok{var(}\VariableTok{{-}{-}basic{-}sidebar{-}logo{-}text}\FunctionTok{)}\OperatorTok{;} \CommentTok{/* 把变量写进伪元素。 */}
  \KeywordTok{writing{-}mode}\CharTok{:} \DecValTok{vertical{-}rl}\OperatorTok{;} \CommentTok{/* 让短文本竖排。 */}
  \KeywordTok{letter{-}spacing}\CharTok{:} \DecValTok{0.18}\DataTypeTok{em}\OperatorTok{;} \CommentTok{/* 增加字间距。 */}
  \KeywordTok{opacity}\CharTok{:} \DecValTok{0.55}\OperatorTok{;} \CommentTok{/* 保持装饰性质。 */}
\NormalTok{\}}
\end{Highlighting}
\end{Shaded}

字体应优先使用系统字体或主题已经提供的短文本字体。徽标文字过长会挤压侧栏。

\subsubsection{12.2 图片徽标}\label{122-ux56feux7247ux5fbdux6807}

\begin{Shaded}
\begin{Highlighting}
\NormalTok{html}\ExtensionTok{[}\SpecialStringTok{data{-}theme}\OperatorTok{=}\StringTok{"advanced"}\ExtensionTok{]}\NormalTok{ \{}
  \VariableTok{{-}{-}advanced{-}sidebar{-}logo{-}image}\CharTok{:} \FunctionTok{url(}\StringTok{"logo.svg"}\FunctionTok{)}\OperatorTok{;} \CommentTok{/* 本地 SVG。 */}
  \VariableTok{{-}{-}advanced{-}sidebar{-}logo{-}text}\CharTok{:} \StringTok{""}\OperatorTok{;}               \CommentTok{/* 禁用备用文字。 */}
\NormalTok{\}}

\NormalTok{html}\ExtensionTok{[}\SpecialStringTok{data{-}theme}\OperatorTok{=}\StringTok{"advanced"}\ExtensionTok{]} \FunctionTok{.glacier{-}sidebar}\InformationTok{::after}\NormalTok{ \{}
  \KeywordTok{content}\CharTok{:} \StringTok{""}\OperatorTok{;} \CommentTok{/* 使用图片时不输出文字。 */}
  \KeywordTok{background}\CharTok{:} \FunctionTok{var(}\VariableTok{{-}{-}advanced{-}sidebar{-}logo{-}image}\FunctionTok{)} \DecValTok{center}\NormalTok{ / }\DecValTok{70}\DataTypeTok{\%} \DecValTok{70}\DataTypeTok{\%} \DecValTok{no{-}repeat}\OperatorTok{;}
  \KeywordTok{opacity}\CharTok{:} \DecValTok{0.9}\OperatorTok{;} \CommentTok{/* 品牌标记可以比普通装饰更清晰。 */}
\NormalTok{\}}
\end{Highlighting}
\end{Shaded}

SVG 应尽量使用简单路径，不携带脚本、外链或嵌入 HTML。复杂插画用
WebP，不要把大图强行做成 SVG。

\begin{center}\rule{0.5\linewidth}{0.5pt}\end{center}

\subsection{13.
图表、曲线和历史页}\label{13-ux56feux8868ux66f2ux7ebfux548cux5386ux53f2ux9875}

曲线页是最容易被装饰破坏可读性的页面。主题作者应先确保坐标轴、曲线、点、悬浮提示和下方控制区域清晰，再考虑背景。

\subsubsection{13.1
曲线编辑器}\label{131-ux66f2ux7ebfux7f16ux8f91ux5668}

\begin{Shaded}
\begin{Highlighting}
\CommentTok{/* 外层只做定位锚点，不绘制一个额外的方形背景。 */}
\NormalTok{html}\ExtensionTok{[}\SpecialStringTok{data{-}theme}\OperatorTok{=}\StringTok{"advanced"}\ExtensionTok{]} \ExtensionTok{[}\SpecialStringTok{data{-}theme{-}card}\OperatorTok{=}\StringTok{"curve{-}editor"}\ExtensionTok{]}\NormalTok{ \{}
  \KeywordTok{background}\CharTok{:} \DecValTok{transparent}\OperatorTok{;}
  \KeywordTok{box{-}shadow}\CharTok{:} \DecValTok{none}\OperatorTok{;}
\NormalTok{\}}

\CommentTok{/* 内层才是真正的图表卡片，需要负责圆角和裁切。 */}
\NormalTok{html}\ExtensionTok{[}\SpecialStringTok{data{-}theme}\OperatorTok{=}\StringTok{"advanced"}\ExtensionTok{]} \ExtensionTok{[}\SpecialStringTok{data{-}theme{-}card}\OperatorTok{=}\StringTok{"curve{-}editor"}\ExtensionTok{]} \OperatorTok{\textgreater{}} \FunctionTok{.relative.rounded{-}3xl}\NormalTok{ \{}
  \KeywordTok{background}\CharTok{:} \FunctionTok{var(}\VariableTok{{-}{-}advanced{-}chart{-}background}\FunctionTok{)}\OperatorTok{;}
  \KeywordTok{border{-}radius}\CharTok{:} \DecValTok{1.25}\DataTypeTok{rem}\OperatorTok{;}
  \KeywordTok{overflow}\CharTok{:} \DecValTok{hidden}\OperatorTok{;}
\NormalTok{\}}
\end{Highlighting}
\end{Shaded}

\subsubsection{13.2 手动档位}\label{132-ux624bux52a8ux6863ux4f4d}

\begin{Shaded}
\begin{Highlighting}
\NormalTok{html}\ExtensionTok{[}\SpecialStringTok{data{-}theme}\OperatorTok{=}\StringTok{"advanced"}\ExtensionTok{]} \ExtensionTok{[}\SpecialStringTok{data{-}theme{-}card}\OperatorTok{=}\StringTok{"curve{-}manual{-}gears"}\ExtensionTok{]}\NormalTok{ \{}
  \KeywordTok{background}\CharTok{:} \FunctionTok{var(}\VariableTok{{-}{-}advanced{-}control{-}background}\FunctionTok{)}\OperatorTok{;}
\NormalTok{\}}

\NormalTok{html}\ExtensionTok{[}\SpecialStringTok{data{-}theme}\OperatorTok{=}\StringTok{"advanced"}\ExtensionTok{]} \ExtensionTok{[}\SpecialStringTok{data{-}theme{-}ui}\OperatorTok{=}\StringTok{"manual{-}gear{-}dot"}\ExtensionTok{]}\NormalTok{ \{}
  \KeywordTok{border{-}color}\CharTok{:} \FunctionTok{var(}\VariableTok{{-}{-}advanced{-}accent}\FunctionTok{)}\OperatorTok{;}
  \KeywordTok{box{-}shadow}\CharTok{:} \DecValTok{0} \DecValTok{0} \DecValTok{0} \DecValTok{3}\DataTypeTok{px} \FunctionTok{rgb(}\DecValTok{47} \DecValTok{109} \DecValTok{246}\NormalTok{ / }\DecValTok{12}\DataTypeTok{\%}\FunctionTok{)}\OperatorTok{;}
\NormalTok{\}}
\end{Highlighting}
\end{Shaded}

手动档位点是操作目标，不能用装饰覆盖，也不要在 hover
时改变点的外部尺寸。

\subsubsection{13.3 历史趋势}\label{133-ux5386ux53f2ux8d8bux52bf}

\begin{Shaded}
\begin{Highlighting}
\NormalTok{html}\ExtensionTok{[}\SpecialStringTok{data{-}theme}\OperatorTok{=}\StringTok{"advanced"}\ExtensionTok{]} \ExtensionTok{[}\SpecialStringTok{data{-}theme{-}card}\OperatorTok{=}\StringTok{"curve{-}history"}\ExtensionTok{]}\NormalTok{ \{}
  \KeywordTok{background}\CharTok{:} \FunctionTok{var(}\VariableTok{{-}{-}advanced{-}history{-}background}\FunctionTok{)}\OperatorTok{;}
\NormalTok{\}}

\NormalTok{html}\ExtensionTok{[}\SpecialStringTok{data{-}theme}\OperatorTok{=}\StringTok{"advanced"}\ExtensionTok{]} \ExtensionTok{[}\SpecialStringTok{data{-}theme{-}ui}\OperatorTok{=}\StringTok{"history{-}display{-}dialog"}\ExtensionTok{]}\NormalTok{ \{}
  \KeywordTok{background}\CharTok{:} \FunctionTok{var(}\VariableTok{{-}{-}advanced{-}dialog{-}background}\FunctionTok{)}\OperatorTok{;}
  \KeywordTok{color}\CharTok{:} \FunctionTok{var(}\VariableTok{{-}{-}foreground}\FunctionTok{)}\OperatorTok{;}
\NormalTok{\}}

\NormalTok{html}\ExtensionTok{[}\SpecialStringTok{data{-}theme}\OperatorTok{=}\StringTok{"advanced"}\ExtensionTok{]} \ExtensionTok{[}\SpecialStringTok{data{-}theme{-}ui}\OperatorTok{=}\StringTok{"history{-}display{-}row"}\ExtensionTok{]}\NormalTok{ \{}
  \KeywordTok{border{-}color}\CharTok{:} \FunctionTok{var(}\VariableTok{{-}{-}advanced{-}border}\FunctionTok{)}\OperatorTok{;}
\NormalTok{\}}
\end{Highlighting}
\end{Shaded}

历史图的功耗、温度和风扇趋势可能分成上下两个图。不要用同一张大背景图压住坐标轴，也不要给不同曲线使用相近颜色。

\begin{center}\rule{0.5\linewidth}{0.5pt}\end{center}

\subsection{14.
设置页和设备连接页}\label{14-ux8bbeux7f6eux9875ux548cux8bbeux5907ux8fdeux63a5ux9875}

设置页是信息密度最高的页面，视觉主题必须优先保护文字、下拉、开关、错误提示和滚动位置。

\subsubsection{14.1 设置分区}\label{141-ux8bbeux7f6eux5206ux533a}

\begin{Shaded}
\begin{Highlighting}
\NormalTok{html}\ExtensionTok{[}\SpecialStringTok{data{-}theme}\OperatorTok{=}\StringTok{"advanced"}\ExtensionTok{]} \ExtensionTok{[}\SpecialStringTok{data{-}theme{-}ui}\OperatorTok{=}\StringTok{"setting{-}section"}\ExtensionTok{]}\NormalTok{ \{}
  \KeywordTok{background}\CharTok{:} \FunctionTok{var(}\VariableTok{{-}{-}advanced{-}setting{-}section{-}background}\FunctionTok{)}\OperatorTok{;}
  \KeywordTok{border{-}color}\CharTok{:} \FunctionTok{var(}\VariableTok{{-}{-}advanced{-}border}\FunctionTok{)}\OperatorTok{;}
  \KeywordTok{box{-}shadow}\CharTok{:} \FunctionTok{var(}\VariableTok{{-}{-}advanced{-}setting{-}shadow}\FunctionTok{)}\OperatorTok{;}
\NormalTok{\}}

\NormalTok{html}\ExtensionTok{[}\SpecialStringTok{data{-}theme}\OperatorTok{=}\StringTok{"advanced"}\ExtensionTok{]} \ExtensionTok{[}\SpecialStringTok{data{-}theme{-}ui}\OperatorTok{=}\StringTok{"setting{-}section{-}header"}\ExtensionTok{]}\NormalTok{ \{}
  \KeywordTok{border{-}color}\CharTok{:} \FunctionTok{var(}\VariableTok{{-}{-}advanced{-}border}\FunctionTok{)}\OperatorTok{;}
  \KeywordTok{color}\CharTok{:} \FunctionTok{var(}\VariableTok{{-}{-}foreground}\FunctionTok{)}\OperatorTok{;}
\NormalTok{\}}
\end{Highlighting}
\end{Shaded}

设置分区的外框应该稳定。不要在切换分区时通过 CSS
动画改变高度，也不要让标题栏在 hover 时产生位移。

\subsubsection{14.2
设备连接面板}\label{142-ux8bbeux5907ux8fdeux63a5ux9762ux677f}

\begin{Shaded}
\begin{Highlighting}
\NormalTok{html}\ExtensionTok{[}\SpecialStringTok{data{-}theme}\OperatorTok{=}\StringTok{"advanced"}\ExtensionTok{]} \ExtensionTok{[}\SpecialStringTok{data{-}theme{-}ui}\OperatorTok{=}\StringTok{"connection{-}panel"}\ExtensionTok{]}\NormalTok{ \{}
  \KeywordTok{background}\CharTok{:} \FunctionTok{var(}\VariableTok{{-}{-}advanced{-}connection{-}background}\FunctionTok{)}\OperatorTok{;}
  \KeywordTok{border{-}color}\CharTok{:} \FunctionTok{var(}\VariableTok{{-}{-}advanced{-}border}\FunctionTok{)}\OperatorTok{;}
\NormalTok{\}}

\NormalTok{html}\ExtensionTok{[}\SpecialStringTok{data{-}theme}\OperatorTok{=}\StringTok{"advanced"}\ExtensionTok{]} \ExtensionTok{[}\SpecialStringTok{data{-}theme{-}ui}\OperatorTok{=}\StringTok{"compatibility{-}submenu"}\ExtensionTok{]}\NormalTok{ \{}
  \KeywordTok{background}\CharTok{:} \FunctionTok{var(}\VariableTok{{-}{-}advanced{-}submenu{-}background}\FunctionTok{)}\OperatorTok{;}
  \KeywordTok{border{-}color}\CharTok{:} \FunctionTok{var(}\VariableTok{{-}{-}advanced{-}border}\FunctionTok{)}\OperatorTok{;}
\NormalTok{\}}
\end{Highlighting}
\end{Shaded}

兼容模式下可能显示
IP、连接状态、扫描进度和重试按钮。装饰层必须位于这些控件之后，不能遮住状态文字。

\begin{center}\rule{0.5\linewidth}{0.5pt}\end{center}

\subsection{15.
颜色、透明度、阴影和动效}\label{15-ux989cux8272ux900fux660eux5ea6ux9634ux5f71ux548cux52a8ux6548}

\subsubsection{15.1 颜色对比}\label{151-ux989cux8272ux5bf9ux6bd4}

至少检查以下组合：

\begin{longtable}[]{@{}L{0.18\textwidth}L{0.30\textwidth}L{0.40\textwidth}@{}}
\toprule\noalign{}\rowcolor{tablehead}
前景 & 背景 & 必须检查 \\
\midrule\noalign{}
\endhead
\bottomrule\noalign{}
\endlastfoot
标题文字 & 首页主卡片 & 设备名和模式名称清晰 \\
次级文字 & 半透明卡片 & 描述和单位清晰 \\
图表线 & 图表背景 & 多条线能区分 \\
开关滑块 & 开关轨道 & 选中态明显 \\
错误文字 & 错误背景 & 不只依靠红色，还要有文字 \\
focus ring & 控件边界 & 键盘操作可见 \\
\end{longtable}

不要把可读性寄托在背景图片的``平均颜色''上。图片会因窗口尺寸、裁切和系统
DPI 改变。

\subsubsection{15.2 阴影}\label{152-ux9634ux5f71}

阴影应表达层级，而不是装饰一切：

\begin{Shaded}
\begin{Highlighting}
\NormalTok{html}\ExtensionTok{[}\SpecialStringTok{data{-}theme}\OperatorTok{=}\StringTok{"advanced"}\ExtensionTok{]} \FunctionTok{.glacier{-}chart{-}card}\NormalTok{ \{}
  \KeywordTok{box{-}shadow}\CharTok{:}
    \DecValTok{0} \DecValTok{8}\DataTypeTok{px} \DecValTok{24}\DataTypeTok{px} \FunctionTok{rgb(}\DecValTok{28} \DecValTok{51} \DecValTok{83}\NormalTok{ / }\DecValTok{10}\DataTypeTok{\%}\FunctionTok{)}\OperatorTok{,} \CommentTok{/* 外部层次。 */}
    \DecValTok{inset} \DecValTok{0} \DecValTok{1}\DataTypeTok{px} \DecValTok{0} \FunctionTok{rgb(}\DecValTok{255} \DecValTok{255} \DecValTok{255}\NormalTok{ / }\DecValTok{45}\DataTypeTok{\%}\FunctionTok{)}\OperatorTok{;} \CommentTok{/* 内部高光。 */}
\NormalTok{\}}
\end{Highlighting}
\end{Shaded}

过重的阴影会让透明背景变成灰块。设置页和曲线页应比首页少使用装饰阴影。

\subsubsection{15.3
动效和减少动态效果}\label{153-ux52a8ux6548ux548cux51cfux5c11ux52a8ux6001ux6548ux679c}

\begin{Shaded}
\begin{Highlighting}
\NormalTok{html}\ExtensionTok{[}\SpecialStringTok{data{-}theme}\OperatorTok{=}\StringTok{"advanced"}\ExtensionTok{]} \ExtensionTok{[}\SpecialStringTok{data{-}theme{-}card}\OperatorTok{=}\StringTok{"fan{-}curve{-}preview"}\ExtensionTok{]}\NormalTok{ \{}
  \KeywordTok{transition}\CharTok{:} \DecValTok{box{-}shadow 180}\DataTypeTok{ms} \DecValTok{ease}\OperatorTok{,} \DecValTok{transform 180}\DataTypeTok{ms} \DecValTok{ease}\OperatorTok{;}
\NormalTok{\}}

\NormalTok{html}\ExtensionTok{[}\SpecialStringTok{data{-}theme}\OperatorTok{=}\StringTok{"advanced"}\ExtensionTok{]} \ExtensionTok{[}\SpecialStringTok{data{-}theme{-}card}\OperatorTok{=}\StringTok{"fan{-}curve{-}preview"}\ExtensionTok{]}\InformationTok{:hover}\NormalTok{ \{}
  \KeywordTok{transform}\CharTok{:} \FunctionTok{translateY(}\DecValTok{{-}1}\DataTypeTok{px}\FunctionTok{)}\OperatorTok{;} \CommentTok{/* 只做轻微位移。 */}
\NormalTok{\}}

\ImportTok{@media} \FunctionTok{(}\KeywordTok{prefers{-}reduced{-}motion}\CharTok{:} \DecValTok{reduce}\FunctionTok{)}\NormalTok{ \{}
\NormalTok{  html}\ExtensionTok{[}\SpecialStringTok{data{-}theme}\OperatorTok{=}\StringTok{"advanced"}\ExtensionTok{]} \OperatorTok{*}\NormalTok{ \{}
    \KeywordTok{animation{-}duration}\CharTok{:} \DecValTok{0.01}\DataTypeTok{ms} \AttributeTok{!important}\OperatorTok{;}
    \KeywordTok{animation{-}iteration{-}count}\CharTok{:} \DecValTok{1} \AttributeTok{!important}\OperatorTok{;}
    \KeywordTok{transition{-}duration}\CharTok{:} \DecValTok{0.01}\DataTypeTok{ms} \AttributeTok{!important}\OperatorTok{;}
\NormalTok{  \}}
\NormalTok{\}}
\end{Highlighting}
\end{Shaded}

不要把扫描光、呼吸动画和大背景动画同时放在每个卡片上。动画应该是可关闭的增强效果，而不是信息传递的唯一方式。

\begin{center}\rule{0.5\linewidth}{0.5pt}\end{center}

\subsection{16. 资源路径、图片和
SVG}\label{16-ux8d44ux6e90ux8defux5f84ux56feux7247ux548c-svg}

\subsubsection{16.1
只使用相对路径}\label{161-ux53eaux4f7fux7528ux76f8ux5bf9ux8defux5f84}

正确：

\begin{Shaded}
\begin{Highlighting}
\NormalTok{html}\ExtensionTok{[}\SpecialStringTok{data{-}theme}\OperatorTok{=}\StringTok{"cute{-}blue"}\ExtensionTok{]}\NormalTok{ \{}
  \VariableTok{{-}{-}cute{-}blue{-}hero}\CharTok{:} \FunctionTok{url(}\StringTok{"hero.webp"}\FunctionTok{)}\OperatorTok{;}
  \VariableTok{{-}{-}cute{-}blue{-}star}\CharTok{:} \FunctionTok{url(}\StringTok{"decorations/star.svg"}\FunctionTok{)}\OperatorTok{;}
  \VariableTok{{-}{-}cute{-}blue{-}font}\CharTok{:} \FunctionTok{url(}\StringTok{"assets/fonts/cute{-}blue{-}ui.woff2"}\FunctionTok{)}\OperatorTok{;}
\NormalTok{\}}
\end{Highlighting}
\end{Shaded}

错误：

\begin{Shaded}
\begin{Highlighting}
\CommentTok{/* 绝对 Windows 路径：只在作者电脑上有效。 */}
\NormalTok{{-}{-}hero}\InformationTok{:}\NormalTok{ url}\InformationTok{(}\NormalTok{"C}\InformationTok{:}\NormalTok{/Users/Alice/Desktop/hero}\FunctionTok{.webp}\NormalTok{"}\InformationTok{)}\NormalTok{;}

\CommentTok{/* 外链：会受网络、域名和隐私策略影响。 */}
\NormalTok{{-}{-}hero}\InformationTok{:}\NormalTok{ url}\InformationTok{(}\NormalTok{"https}\InformationTok{:}\ErrorTok{//example.com/hero.webp");}
\end{Highlighting}
\end{Shaded}

\subsubsection{16.2 WebP 和 SVG
的选择}\label{162-webp-ux548c-svg-ux7684ux9009ux62e9}

\begin{longtable}[]{@{}L{0.18\textwidth}L{0.30\textwidth}L{0.40\textwidth}@{}}
\toprule\noalign{}\rowcolor{tablehead}
资源 & 推荐格式 & 说明 \\
\midrule\noalign{}
\endhead
\bottomrule\noalign{}
\endlastfoot
复杂插画、照片、纹理 & WebP & 文件小，适合复杂像素内容 \\
线稿、简单徽标、星星和云朵 & SVG & 可缩放，适合少量路径 \\
字体 & WOFF2 & 体积小，浏览器支持好 \\
透明复杂贴纸 & WebP 或 SVG & 视细节数量选择 \\
\end{longtable}

\subsubsection{16.3
资源大小建议}\label{163-ux8d44ux6e90ux5927ux5c0fux5efaux8bae}

分享包和便携包都应该优先使用小资源：

\begin{longtable}[]{@{}L{0.36\textwidth}L{0.58\textwidth}@{}}
\toprule\noalign{}\rowcolor{tablehead}
资源 & 建议上限 \\
\midrule\noalign{}
\endhead
\bottomrule\noalign{}
\endlastfoot
\texttt{hero.webp} & 约 500 KB \\
单个贴纸 & 约 200 KB \\
单个 SVG & 约 100 KB，尽量更小 \\
单个字体子集 & 约 300 KB \\
完整主题包 & 建议控制在 5 MB 内 \\
\end{longtable}

这些是工程建议，不是软件强制的二进制限制；自动校验可以根据实际发布渠道调整。

\begin{center}\rule{0.5\linewidth}{0.5pt}\end{center}

\subsection{17.
字体和字体子集}\label{17-ux5b57ux4f53ux548cux5b57ux4f53ux5b50ux96c6}

\subsubsection{17.1 推荐顺序}\label{171-ux63a8ux8350ux987aux5e8f}

\begin{enumerate}
\def\labelenumi{\arabic{enumi}.}
\tightlist
\item
  系统字体：体积为零，跨主题最稳定。
\item
  小型 WOFF2 子集：用于标题、主题名或少量装饰文字。
\item
  完整字体：除非许可证和包体积都明确允许，否则不推荐。
\end{enumerate}

\subsubsection{\texorpdfstring{17.2 \texttt{@font-face}
逐行示例}{17.2 @font-face 逐行示例}}\label{172-font-face-ux9010ux884cux793aux4f8b}

\begin{Shaded}
\begin{Highlighting}
\ImportTok{@font{-}face}\NormalTok{ \{}
  \KeywordTok{font{-}family}\CharTok{:} \StringTok{"Cute Blue UI"}\OperatorTok{;} \CommentTok{/* CSS 中引用的字体族名称。 */}
  \KeywordTok{src}\CharTok{:} \FunctionTok{url(}\StringTok{"assets/fonts/cute{-}blue{-}ui.woff2"}\FunctionTok{)} \FunctionTok{format(}\StringTok{"woff2"}\FunctionTok{)}\OperatorTok{;} \CommentTok{/* 本地相对字体。 */}
  \KeywordTok{font{-}weight}\CharTok{:} \DecValTok{400} \DecValTok{700}\OperatorTok{;} \CommentTok{/* 支持普通到粗体的可变范围。 */}
  \KeywordTok{font{-}style}\CharTok{:} \DecValTok{normal}\OperatorTok{;} \CommentTok{/* 不声明斜体，避免浏览器合成假斜体。 */}
  \KeywordTok{font{-}display}\CharTok{:} \DecValTok{swap}\OperatorTok{;} \CommentTok{/* 字体未加载时先用回退字体。 */}
\NormalTok{\}}

\NormalTok{html}\ExtensionTok{[}\SpecialStringTok{data{-}theme}\OperatorTok{=}\StringTok{"cute{-}blue"}\ExtensionTok{]}\NormalTok{ \{}
  \VariableTok{{-}{-}cute{-}blue{-}ui{-}font}\CharTok{:} \StringTok{"Cute Blue UI"}\OperatorTok{,} \StringTok{"Segoe UI"}\OperatorTok{,} \StringTok{"Microsoft YaHei UI"}\OperatorTok{,} \DecValTok{sans{-}serif}\OperatorTok{;}
\NormalTok{\}}

\NormalTok{html}\ExtensionTok{[}\SpecialStringTok{data{-}theme}\OperatorTok{=}\StringTok{"cute{-}blue"}\ExtensionTok{]} \FunctionTok{.glacier{-}titlebar}\OperatorTok{,}
\NormalTok{html}\ExtensionTok{[}\SpecialStringTok{data{-}theme}\OperatorTok{=}\StringTok{"cute{-}blue"}\ExtensionTok{]} \ExtensionTok{[}\SpecialStringTok{data{-}theme{-}ui}\OperatorTok{=}\StringTok{"setting{-}section{-}header"}\ExtensionTok{]}\NormalTok{ \{}
  \KeywordTok{font{-}family}\CharTok{:} \FunctionTok{var(}\VariableTok{{-}{-}cute{-}blue{-}ui{-}font}\FunctionTok{)}\OperatorTok{;} \CommentTok{/* 只把字体用于适合的位置。 */}
\NormalTok{\}}
\end{Highlighting}
\end{Shaded}

正文、设置描述和错误提示建议保留系统回退字体。主题字体只覆盖标题或明确的短文本区域，避免字形缺失导致整页方块。

\subsubsection{17.3
字体子集脚本的设计}\label{173-ux5b57ux4f53ux5b50ux96c6ux811aux672cux7684ux8bbeux8ba1}

仓库可以提供 \texttt{tools/subset\_theme\_font.py}，开发者本地安装
\texttt{fonttools} 后运行。脚本不进入 FanControl
运行时，也不会增加安装包体积。

建议输入：

\begin{Shaded}
\begin{Highlighting}
\NormalTok{原始字体文件 + 主题目录 + 主题模式(title/ui) + 可选文本文件}
\end{Highlighting}
\end{Shaded}

建议输出：

\begin{Shaded}
\begin{Highlighting}
\NormalTok{assets/fonts/\textless{}theme{-}id\textgreater{}{-}\textless{}mode\textgreater{}{-}subset.woff2}
\NormalTok{assets/fonts/LICENSE{-}\textless{}font{-}name\textgreater{}.txt}
\end{Highlighting}
\end{Shaded}

脚本应完成以下工作：

\begin{enumerate}
\def\labelenumi{\arabic{enumi}.}
\tightlist
\item
  读取 \texttt{theme.json} 的 \texttt{name}、\texttt{description}
  和作者说明。
\item
  扫描 CSS 中的 \texttt{content:\ "..."} 字符和显式标题文本。
\item
  合并 ASCII、数字、中文标点、主题专用字符和作者提供的文本文件。
\item
  调用 \texttt{fontTools.pyftsubset}，输出 WOFF2。
\item
  生成缺字报告，列出无法从源字体找到的字符。
\item
  复制或检查字体许可证。
\end{enumerate}

示例命令：

\begin{Shaded}
\begin{Highlighting}
\CommentTok{\# {-}{-}font 是原始字体，{-}{-}theme 是主题目录，{-}{-}mode 决定字符集规模。}
\NormalTok{python tools}\OperatorTok{/}\NormalTok{subset\_theme\_font}\OperatorTok{.}\FunctionTok{py}\NormalTok{ \textasciigrave{}}
  \OperatorTok{{-}{-}}\NormalTok{font source}\OperatorTok{.}\FunctionTok{ttf}\NormalTok{ \textasciigrave{}}
  \OperatorTok{{-}{-}}\NormalTok{theme themes}\OperatorTok{/}\NormalTok{cute{-}blue \textasciigrave{}}
  \OperatorTok{{-}{-}}\NormalTok{mode ui \textasciigrave{}}
  \OperatorTok{{-}{-}}\NormalTok{output themes}\OperatorTok{/}\NormalTok{cute{-}blue}\OperatorTok{/}\NormalTok{assets}\OperatorTok{/}\NormalTok{fonts}\OperatorTok{/}\NormalTok{cute{-}blue{-}ui{-}subset}\OperatorTok{.}\FunctionTok{woff2}
\end{Highlighting}
\end{Shaded}

\texttt{title} 模式只保留主题名和装饰标题，包体积最小；\texttt{ui}
模式应额外包含软件常用标签和作者提供的界面文案。中文 UI
主题必须考虑用户切换语言后的英文、日文和标点回退。

\subsubsection{17.4
字体缺字排查}\label{174-ux5b57ux4f53ux7f3aux5b57ux6392ux67e5}

\begin{longtable}[]{@{}L{0.18\textwidth}L{0.30\textwidth}L{0.40\textwidth}@{}}
\toprule\noalign{}\rowcolor{tablehead}
现象 & 原因 & 修复 \\
\midrule\noalign{}
\endhead
\bottomrule\noalign{}
\endlastfoot
主题名显示方块 & 子集没有主题名字符 & 把名称加入文本集合后重新裁剪 \\
标题正常、正文全变系统字体 & 字体只覆盖标题或正文缺字 & 扩大 UI
子集或缩小字体作用范围 \\
加粗时字形变化异常 & 子集没有对应权重或变量轴 & 保留正确
\texttt{font-weight} 范围 \\
首次打开字体闪烁 & 字体尚未加载 & 保留 \texttt{font-display:\ swap}
和系统回退 \\
包体积突然增大 & 打包了完整字体 & 使用 \texttt{pyftsubset} 生成 WOFF2 \\
\end{longtable}

\begin{center}\rule{0.5\linewidth}{0.5pt}\end{center}

\subsection{18.
外部资源和安全边界}\label{18-ux5916ux90e8ux8d44ux6e90ux548cux5b89ux5168ux8fb9ux754c}

主题 CSS 会被注入 GUI 页面，因此``只是 CSS''不等于完全没有风险。

\subsubsection{18.1
可分享主题应禁止}\label{181-ux53efux5206ux4eabux4e3bux9898ux5e94ux7981ux6b62}

\begin{itemize}
\tightlist
\item
  \texttt{@import} 外部样式表
\item
  \texttt{url(http://...)}、\texttt{url(https://...)} 和远程字体
\item
  JavaScript、HTML、可执行文件和 DLL
\item
  绝对文件路径
\item
  符号链接和 ZIP 路径穿越
\item
  未声明许可证的第三方字体和图片
\item
  覆盖官方主题 ID 的非官方包
\end{itemize}

本地 DIY 可以暂时保留实验性写法，但准备分享主题时必须改成本地资源。

\subsubsection{18.2 为什么不允许远程
URL}\label{182-ux4e3aux4ec0ux4e48ux4e0dux5141ux8bb8ux8fdcux7a0b-url}

远程 URL
会带来网络不可用、第三方服务下线、隐私请求、加载延迟和主题不可复现问题。主题运行时资源应当随主题包提供，并通过本地资源路由读取。

\subsubsection{18.3 ZIP
安装安全}\label{183-zip-ux5b89ux88c5ux5b89ux5168}

任何未来的安装器或手动解压流程都应该检查：

\begin{enumerate}
\def\labelenumi{\arabic{enumi}.}
\tightlist
\item
  每个路径都是相对路径。
\item
  清理后的路径仍然位于临时主题目录内。
\item
  不接受符号链接和特殊文件。
\item
  文件总数、单文件大小和解压后总大小有上限。
\item
  必须存在 \texttt{theme.json} 和 \texttt{theme.css}。
\item
  安装完成后使用临时目录原子替换，失败时保留旧主题。
\end{enumerate}

\begin{center}\rule{0.5\linewidth}{0.5pt}\end{center}

\subsection{19.
当前产品边界}\label{19-ux5f53ux524dux4ea7ux54c1ux8fb9ux754c}

本文暂不规定主题市场、联网清单或上传流程。主题作者当前只需要准备符合本文规范的主题目录或
ZIP；任何发布渠道都应以本指南的文件、资源和安全约束为准。

\begin{center}\rule{0.5\linewidth}{0.5pt}\end{center}

\subsection{20.
本地调试方法}\label{20-ux672cux5730ux8c03ux8bd5ux65b9ux6cd5}

\subsubsection{21.1
让主题出现在列表}\label{211-ux8ba9ux4e3bux9898ux51faux73b0ux5728ux5217ux8868}

确认主题目录位于可执行文件同级：

\begin{Shaded}
\begin{Highlighting}
\NormalTok{FanControl/}
\NormalTok{  FanControl.exe}
\NormalTok{  themes/}
\NormalTok{    cute{-}blue/}
\NormalTok{      theme.json}
\NormalTok{      theme.css}
\end{Highlighting}
\end{Shaded}

启动后打开设置页的主题选择器。如果主题没有出现，先检查 JSON
是否能被解析，再检查 ID 是否符合规则。

\subsubsection{21.2 CSS 不生效}\label{212-css-ux4e0dux751fux6548}

按以下顺序检查：

\begin{enumerate}
\def\labelenumi{\arabic{enumi}.}
\tightlist
\item
  文件名是否严格为 \texttt{theme.css}。
\item
  \texttt{theme.json} 的 \texttt{id} 是否与文件夹相同。
\item
  CSS 是否使用相同的 \texttt{html{[}data-theme="id"{]}} 前缀。
\item
  自定义主题是否真正被选中，而不是仍然是 \texttt{system}、\texttt{light}
  或 \texttt{dark}。
\item
  浏览器开发者工具中是否存在 \texttt{html{[}data-theme="id"{]}}。
\item
  是否被模板文件末尾的最终覆盖规则覆盖。
\end{enumerate}

\subsubsection{21.3
图片或字体不显示}\label{213-ux56feux7247ux6216ux5b57ux4f53ux4e0dux663eux793a}

检查 CSS 的相对路径和真实文件大小写：

\begin{Shaded}
\begin{Highlighting}
\NormalTok{theme.css                         \# 引用 url("hero.webp")}
\NormalTok{hero.webp                         \# 必须位于同一个主题目录}
\NormalTok{assets/fonts/ui.woff2             \# 路径必须与 CSS 完全一致}
\end{Highlighting}
\end{Shaded}

不要把 \texttt{/theme-assets/...} 手动写成其他主题
ID。应用会把同主题的相对路径自动改写为资源路由。

\subsubsection{21.4
圆角露底或出现方角}\label{214-ux5706ux89d2ux9732ux5e95ux6216ux51faux73b0ux65b9ux89d2}

先确认装饰元素和外壳的层级，再检查真实绘制卡片。曲线页常见结构是外层
\texttt{data-theme-card} 属性（值为 \texttt{curve-editor}）加内部圆角图表卡片，
因此只修改外层不一定能解决问题。

\subsubsection{21.5 页面变卡}\label{215-ux9875ux9762ux53d8ux5361}

优先减少：

\begin{itemize}
\tightlist
\item
  大图尺寸和数量
\item
  \texttt{backdrop-filter:\ blur(...)}
\item
  持续运行的 \texttt{@keyframes}
\item
  每个卡片都使用的伪元素
\item
  多层 box-shadow
\end{itemize}

不要先重写整个页面或把图表改成另一种渲染技术。主题 CSS
的目标是视觉覆盖，不是重构应用布局。

\begin{center}\rule{0.5\linewidth}{0.5pt}\end{center}

\subsection{21.
代码注释逐类说明}\label{21-ux4ee3ux7801ux6ce8ux91caux9010ux7c7bux8bf4ux660e}

模板中的注释不是装饰，它们对应了维护时最容易出问题的区域。新主题应继续保留这种分区注释。

\subsubsection{22.1 Basic
模板注释分组}\label{221-basic-ux6a21ux677fux6ce8ux91caux5206ux7ec4}

\begin{longtable}[]{@{}L{0.36\textwidth}L{0.58\textwidth}@{}}
\toprule\noalign{}\rowcolor{tablehead}
注释主题 & 代码意思 \\
\midrule\noalign{}
\endhead
\bottomrule\noalign{}
\endlastfoot
文件头和主题 ID & 提醒复制主题时同步修改 JSON ID 和 CSS 前缀 \\
滚动条 & 说明滚动条不应使用过重颜色抢走视觉 \\
图表颜色 & 说明 CPU、GPU、风扇和功耗曲线需要区分 \\
次级背景和弹窗 & 说明下拉、toast 和辅助按钮使用的层级 \\
左侧导航栏 & 说明侧栏颜色和选中状态 \\
内容面板和卡片 & 说明基础版尽量不依赖重 blur \\
层级变量 & 提醒不要让装饰盖住弹窗和控件 \\
图标描边 & 说明描边粗细只能调视觉，不能改变图标盒子尺寸 \\
页面背景 & 说明页面属性在旧版本中可能被忽略 \\
控件 hover/focus & 说明 hover 不应造成布局移动，focus 必须可见 \\
开关兼容写法 & 说明 Basic 使用 \texttt{button{[}role="switch"{]}}
兼容旧版本 \\
左侧标记 & 说明文字徽标和图片徽标的二选一关系 \\
仪表和半环 & 说明只改变颜色和光效，不重建 SVG 或图表结构 \\
\end{longtable}

\subsubsection{22.2 Advanced
模板注释分组}\label{222-advanced-ux6a21ux677fux6ce8ux91caux5206ux7ec4}

\begin{longtable}[]{@{}L{0.36\textwidth}L{0.58\textwidth}@{}}
\toprule\noalign{}\rowcolor{tablehead}
注释主题 & 代码意思 \\
\midrule\noalign{}
\endhead
\bottomrule\noalign{}
\endlastfoot
外壳圆角填充 & 防止透明窗口在四角露出原生背景 \\
高级外观层 & 集中放贴纸、幕布和玻璃效果，方便整段删除 \\
字体与卡片微调 & 控制 UI 字体、透明度、hover 稳定性 \\
字体子集 & 说明 WOFF2 子集的目的和完整字体的体积风险 \\
透明卡片层 & 让背景透出但保留文字可读性 \\
幕布可见性 & 说明觉得太花时可以删除整段 \\
左侧按钮点击反馈 & 说明轻微缩放和光圈不应改变布局 \\
曲线编辑器修复 & 说明外层是布局锚点，内部才是真实圆角卡片 \\
\end{longtable}

\subsubsection{22.3
注释的写法}\label{223-ux6ce8ux91caux7684ux5199ux6cd5}

推荐写``原因 + 约束''，不推荐写无信息量的句子：

\begin{Shaded}
\begin{Highlighting}
\CommentTok{/* 好：伪元素只做装饰，不能挡住按钮点击。 */}
\NormalTok{pointer{-}events}\InformationTok{:}\NormalTok{ none;}

\CommentTok{/* 差：设置 pointer{-}events。 */}
\NormalTok{pointer{-}events}\InformationTok{:}\NormalTok{ none;}
\end{Highlighting}
\end{Shaded}

复杂区块前保留一段短注释即可。不要在每一行都重复 CSS
属性的字面意思；但对于
z-index、透明度、固定尺寸、资源路径和最终覆盖顺序，应该说明为什么这样写。

\begin{center}\rule{0.5\linewidth}{0.5pt}\end{center}

\subsection{22.
完整验收清单}\label{22-ux5b8cux6574ux9a8cux6536ux6e05ux5355}

\subsubsection{23.1
文件和清单}\label{231-ux6587ux4ef6ux548cux6e05ux5355}

\begin{itemize}
\tightlist
\item[$\square$]
  文件夹名与 \texttt{theme.json} 的 \texttt{id} 完全一致。
\item[$\square$]
  \texttt{id} 只包含小写字母、数字、短横线或下划线。
\item[$\square$]
  \texttt{name}、\texttt{author}、\texttt{version}、\texttt{description}
  已填写。
\item[$\square$]
  \texttt{base} 是 \texttt{light} 或 \texttt{dark}。
\item[$\square$]
  \texttt{layer} 是 \texttt{basic} 或 \texttt{advanced}。
\item[$\square$]
  没有在 JSON 中写 \texttt{//} 注释。
\item[$\square$]
  主题包至少包含 \texttt{theme.json} 和 \texttt{theme.css}。
\end{itemize}

\subsubsection{23.2 CSS 和组件}\label{232-css-ux548cux7ec4ux4ef6}

\begin{itemize}
\tightlist
\item[$\square$]
  所有主题选择器都从 \texttt{html{[}data-theme="主题ID"{]}} 开始。
\item[$\square$]
  没有无作用域的 \texttt{body}、\texttt{:root} 或
  \texttt{.glacier-shell} 覆盖。
\item[$\square$]
  页面背景没有遮挡正文。
\item[$\square$]
  卡片透明度足以读取标题、单位和数值。
\item[$\square$]
  下拉框、开关、滑块、按钮和 focus ring 可见。
\item[$\square$]
  设置行 hover 不改变布局尺寸。
\item[$\square$]
  曲线编辑器、历史图和功耗图没有露出方角。
\item[$\square$]
  曲线颜色彼此区分。
\item[$\square$]
  主题在首页、曲线页、设置页和关于页都测试过。
\end{itemize}

\subsubsection{23.3
资源和字体}\label{233-ux8d44ux6e90ux548cux5b57ux4f53}

\begin{itemize}
\tightlist
\item[$\square$]
  所有 \texttt{url(...)} 都是主题包内的相对路径。
\item[$\square$]
  没有绝对路径或远程资源。
\item[$\square$]
  SVG 不含脚本和外链。
\item[$\square$]
  图片经过裁剪和压缩，适合窗口比例。
\item[$\square$]
  字体是授权可分发的字体。
\item[$\square$]
  字体子集包含主题名和主题专用字符。
\item[$\square$]
  字体有系统回退，缺字时不会出现方块。
\item[$\square$]
  \texttt{font-display:\ swap} 已设置。
\end{itemize}

\subsubsection{23.4
性能和可访问性}\label{234-ux6027ux80fdux548cux53efux8bbfux95eeux6027}

\begin{itemize}
\tightlist
\item[$\square$]
  没有在每张卡片上运行持续动画。
\item[$\square$]
  \texttt{backdrop-filter} 使用数量有限。
\item[$\square$]
  \texttt{prefers-reduced-motion} 下动画会关闭或缩短。
\item[$\square$]
  信息不依赖颜色单独表达。
\item[$\square$]
  键盘 focus 状态清楚。
\item[$\square$]
  窄窗口下装饰会隐藏或自动收敛。
\end{itemize}

\subsubsection{23.5
分享前检查}\label{235-ux5206ux4eabux524dux68c0ux67e5}

\begin{itemize}
\tightlist
\item[$\square$]
  预览图是清晰的 WebP，比例稳定。
\item[$\square$]
  包含许可证和作者署名。
\item[$\square$]
  ZIP 内没有开发缓存、工程文件或可执行文件。
\item[$\square$]
  如果发布渠道需要校验，已经生成并记录 SHA-256。
\item[$\square$]
  \texttt{minAppVersion} 与使用的高级钩子匹配。
\item[$\square$]
  更新说明写明视觉变化和兼容性变化。
\item[$\square$]
  提交前在干净的便携目录中安装测试。
\end{itemize}

\begin{center}\rule{0.5\linewidth}{0.5pt}\end{center}

\subsection{23. 给 AI
或协作者的标准提示词}\label{23-ux7ed9-ai-ux6216ux534fux4f5cux8005ux7684ux6807ux51c6ux63d0ux793aux8bcd}

\begin{Shaded}
\begin{Highlighting}
\NormalTok{请为 FanControl 制作一个自定义主题，只修改 theme.json、theme.css 和主题资源，不修改软件源码。}

\NormalTok{请先选择模板：}
\NormalTok{{-} basic：只使用基础颜色变量、glacier{-}* 外壳和 button[role="switch"]，尽量兼容 THRM。}
\NormalTok{{-} advanced：可以使用 data{-}theme{-}page、data{-}theme{-}card、data{-}theme{-}ui、相对图片、贴纸和字体子集。}

\NormalTok{必须遵守：}
\NormalTok{1. 文件夹名、theme.json 的 id 和 html[data{-}theme="id"] 前缀一致。}
\NormalTok{2. 所有 CSS 选择器都放在 html[data{-}theme="id"] 作用域内。}
\NormalTok{3. JSON 不能使用 // 注释，说明请写在 $comment 或 $help。}
\NormalTok{4. 资源只能使用主题包内的相对路径，不能使用 C:/、D:/ 或远程 URL。}
\NormalTok{5. 不要让装饰覆盖按钮、下拉框、开关、进度条、图表和正文。}
\NormalTok{6. 设置页、曲线页、首页和关于页都要保证文字可读。}
\NormalTok{7. 字体优先使用系统字体；使用自定义字体时必须使用授权的 WOFF2 子集。}
\NormalTok{8. 任何 blur、动画、大图和多层阴影都要考虑性能和减少动态效果。}
\NormalTok{9. 每个复杂 CSS 区块前写一条说明“它解决什么问题、为什么放在这里”。}
\NormalTok{10. 最后输出文件清单、许可证、已知兼容版本和验收结果。}

\NormalTok{主题需求：}
\NormalTok{{-} 主题名：}
\NormalTok{{-} 主题 ID：}
\NormalTok{{-} basic 或 advanced：}
\NormalTok{{-} 基础明暗：}
\NormalTok{{-} 主色和辅助色：}
\NormalTok{{-} 首页插画或背景：}
\NormalTok{{-} 左侧徽标：文字还是图片：}
\NormalTok{{-} 是否使用字体子集：}
\NormalTok{{-} 希望强调的页面：}
\NormalTok{{-} 不希望出现的元素：}
\end{Highlighting}
\end{Shaded}

\begin{center}\rule{0.5\linewidth}{0.5pt}\end{center}

\subsection{24. 参考文件}\label{24-ux53c2ux8003ux6587ux4ef6}

\begin{itemize}
\tightlist
\item
  基础模板：\texttt{docs/theme-template/basic/}
\item
  高级模板：\texttt{docs/theme-template/advanced/}
\item
  主题目录：\texttt{themes/}
\item
  已导出主题：\texttt{exported-themes/}
\item
  主题资源和路径实现：\texttt{internal/theme/theme.go}
\item
  主题 GUI API：\texttt{internal/guiapp/theme\_api.go}
\item
  前端主题同步：\texttt{frontend/src/app/components/SystemThemeSync.tsx}
\item
  设置页主题选择：\texttt{frontend/src/app/components/settings/SystemSettingsSection.tsx}
\item
  主题设计指南 PDF 源码：\texttt{docs/theme-diy-guide-latex/}
\end{itemize}

这份指南解释的是当前版本公开的主题接口。新增页面或组件时，应同时更新组件索引、模板示例和验收清单，避免作者只能通过阅读源码猜测钩子含义。

\end{document}
